\newif\ifteacher\teacherfalse
\documentclass[11pt,a4paper]{article}

\usepackage[margin=2.3cm]{geometry}
\usepackage{amsmath,amssymb,amsthm}
\usepackage[T1]{fontenc}
\usepackage{tgschola}
\usepackage{fouriernc}
\usepackage{microtype}
\linespread{1.04}
\usepackage{xcolor}
\usepackage{titlesec}
\usepackage{tikz}
\usetikzlibrary{angles,quotes,calc,arrows.meta}
\usepackage{enumitem}
\usepackage{needspace}
\usepackage[most]{tcolorbox}
\usepackage[hidelinks]{hyperref}

\definecolor{accent}{RGB}{38,52,70}
\definecolor{panel}{RGB}{236,238,241}
\definecolor{muted}{RGB}{70,72,76}

\newtheoremstyle{branddef}{9pt}{9pt}{}{}{\bfseries\color{accent}}{.}{ }{}
\theoremstyle{branddef}
\newtheorem{definition}{Definition}
\newtheorem{example}{Example}

% Theorems: framed, with the NCERT number. #1 = NCERT number, #2 = name of the criterion.
\newenvironment{theorem}[2]
  {\par\medskip\Needspace{5\baselineskip}%
   \begin{tcolorbox}[enhanced,sharp corners,colback=white,colframe=black,boxrule=0.9pt,
     left=7pt,right=7pt,top=5pt,bottom=5pt]%
   \textbf{Theorem #1}\ (#2).\ \itshape}
  {\end{tcolorbox}\par\medskip}
% corollary: same frame, #1 = its name
\newenvironment{corollary}[1]
  {\par\medskip\Needspace{4\baselineskip}%
   \begin{tcolorbox}[enhanced,sharp corners,colback=white,colframe=black,boxrule=0.9pt,
     left=7pt,right=7pt,top=5pt,bottom=5pt]%
   \textbf{Corollary}\ (#1).\ \itshape}
  {\end{tcolorbox}\par\medskip}
\newcommand{\tr}{\triangle}
\newenvironment{exercise}[1]
  {\par\medskip\Needspace{7\baselineskip}\noindent{\large\bfseries Exercise #1}
   \par\smallskip
   \begin{enumerate}[leftmargin=*,itemsep=4pt]}
  {\end{enumerate}}
% Solve-first problems: numbered P1, P2, ... through the whole document
\newcounter{prob}
\newenvironment{solvefirst}
  {\par\smallskip\noindent\textbf{Solve first.}
   \begin{enumerate}[label=\textbf{P\arabic*.},leftmargin=2.2em,itemsep=3pt,start=\numexpr\value{prob}+1\relax]}
  {\setcounter{prob}{\value{enumi}}\end{enumerate}}
\newcommand{\why}{\par\smallskip\noindent\textbf{Why it works.}\par}
\ifteacher
\newenvironment{answerblock}
  {\par\smallskip\Needspace{10\baselineskip}\begingroup\small\linespread{1.15}\selectfont\color{muted}\noindent\textit{Answers.}\par}
  {\par\endgroup\medskip}
\newcommand{\ans}[1]{\par\smallskip\Needspace{9\baselineskip}\noindent\textbf{#1}\ }
\newenvironment{solution}
  {\par\smallskip\begingroup\small\linespread{1.15}\selectfont\color{muted}\noindent\textit{Solution.}\par}
  {\par\endgroup\smallskip}
\else
\usepackage{comment}
\excludecomment{answerblock}
\excludecomment{solution}
\fi
\newcommand{\figph}[1]{\begin{center}\fbox{\parbox{.85\linewidth}{\small FIGURE: #1}}\end{center}}
% right-angle mark at vertex #1 between the directions to #2 and #3
\newcommand{\ramark}[3]{\coordinate (ra1) at ($(#1)!2mm!(#2)$);\coordinate (ra2) at ($(#1)!2mm!(#3)$);\draw[thin] (ra1) -- ($(ra1)+(ra2)-(#1)$) -- (ra2);}
% tick mark across the midpoint of segment #1--#2 (equal lengths)
\newcommand{\tick}[2]{\draw[thin] ($($(#1)!0.5!(#2)$)!1.2mm!90:(#2)$) -- ($($(#1)!0.5!(#2)$)!1.2mm!-90:(#2)$);}
\newcommand{\exfigs}[2]{\begin{center}\begin{tikzpicture}[scale=#1, every node/.style={font=\small}]#2\end{tikzpicture}\end{center}}
\newcommand{\exfig}[1]{\begin{center}\begin{tikzpicture}[scale=.75, every node/.style={font=\small}]#1\end{tikzpicture}\end{center}}
% rough answer diagram (teacher copy only; it sits inside answerblock)
\newcommand{\ansfig}[1]{\begin{center}\begin{tikzpicture}[scale=.6, every node/.style={font=\footnotesize}, draw=muted]#1\end{tikzpicture}\end{center}}
% worked answer (teacher copy): solution text on the left, rough figure on the right
\newcommand{\anspair}[2]{\par\smallskip\noindent
  \begin{minipage}[t]{0.62\linewidth}\vspace{0pt}#1\end{minipage}\hfill
  \begin{minipage}[t]{0.34\linewidth}\vspace{0pt}\centering
  \begin{tikzpicture}[scale=.6, every node/.style={font=\footnotesize}, draw=muted]#2\end{tikzpicture}\end{minipage}\par\smallskip}
\newenvironment{steps}{\begin{enumerate}[label=\arabic*.,leftmargin=1.5em,itemsep=0pt,topsep=2pt]}{\end{enumerate}}
\newcommand{\rememberbox}[1]{%
  \par\medskip\noindent\textcolor{black!35}{\rule{\linewidth}{0.4pt}}\par\nointerlineskip\vspace{4pt}
  \noindent\textbf{Things to Remember}\par\smallskip #1\par\vspace{3pt}
  \noindent\textcolor{black!35}{\rule{\linewidth}{0.4pt}}\par\medskip}
\ifteacher
\newcommand{\facilitatornote}[1]{%
  \par\medskip\noindent\textcolor{black!35}{\rule{\linewidth}{0.4pt}}\par\nointerlineskip\vspace{3pt}
  \noindent{\small\textbf{Note for the teacher.}\ #1}\par\vspace{2pt}
  \noindent\textcolor{black!35}{\rule{\linewidth}{0.4pt}}\par\medskip}
\else
\newcommand{\facilitatornote}[1]{}
\fi

\pagestyle{plain}
\setlength{\parskip}{4pt}
\setlength{\parindent}{0pt}
% keep tall fractions in running text from touching the line above or below
\setlength{\lineskiplimit}{2pt}
\setlength{\lineskip}{3pt}

\titleformat{\section}{\large\bfseries}{\thesection}{0.8em}{}
\titlespacing*{\section}{0pt}{16pt}{6pt}
\newcommand{\sheethead}[1]{\noindent{\LARGE Similar Triangles}\hfill{\large #1}\par\vspace{3pt}
  \noindent\rule{\linewidth}{0.6pt}\par\medskip}
\begin{document}
\sheethead{Class 3}

\ifteacher
\facilitatornote{%
\begin{itemize}[leftmargin=1.4em,itemsep=1pt,topsep=2pt]
  \item \textbf{Aims.} Students write a similarity statement with the vertices in the correct order. They find, by drawing and measuring, that one set of conditions is enough to make two triangles similar: two pairs of equal angles (AA), three pairs of proportional sides (SSS), or one equal angle between two pairs of proportional sides (SAS). They prove the three criteria from the Basic Proportionality Theorem and use them to find unknown angles and lengths, to prove results, and to measure heights from shadows.
  \item \textbf{NCERT links.} NCERT Section 6.4 and Section 6.5. Theorems 6.3, 6.4 and 6.5 keep their NCERT numbers. Examples 1 to 7 of this class are NCERT Examples 4 to 8, together with two examples of our own. The exercises contain all 16 questions of NCERT Exercise 6.3, marked (NCERT 6.3, Q$n$) in this copy only.
  \item \textbf{Prior work.} The meaning of similar figures, the Basic Proportionality Theorem (BPT) and its converse, and the congruence criteria SAS, ASA, SSS and RHS from Class IX.
  \item \textbf{Materials.} A4 sheets, a ruler, a compass and a protractor for each student.
  \item \textbf{Method.} Students solve each \emph{Solve first} set before the theorem is stated. The proofs follow the measurements.
  \item \textbf{Break points.} The material is longer than one class. Natural break points are the end of Exercise 3A (AA criterion), the end of Exercise 3B (SSS criterion) and the end of Exercise 3C (SAS criterion).
\end{itemize}}
\fi

\section{Writing a similarity statement}

Two triangles are \textbf{similar} if (i) their corresponding angles are equal and (ii) their corresponding sides are in the same ratio. We write $\tr ABC\sim\tr DEF$ and read it as ``triangle $ABC$ is similar to triangle $DEF$''. The order of the letters tells us which vertices match: the first letters match, the second letters match and the third letters match.

\begin{solvefirst}
  \item It is given that $\tr ABC\sim\tr PQR$.
  \begin{enumerate}[label=(\alph*),itemsep=1pt]
    \item Write the three pairs of equal angles.
    \item Write the three equal ratios of sides.
    \item Raju writes $\tr ABC\sim\tr QPR$ for the same two triangles. Is he correct? Give a reason.
    \item Is it correct to write $\tr BCA\sim\tr QRP$? Give a reason.
  \end{enumerate}
\begin{solution}
\begin{steps}
  \item[(a)] $\angle A=\angle P$, $\angle B=\angle Q$, $\angle C=\angle R$.
  \item[(b)] $\dfrac{AB}{PQ}=\dfrac{BC}{QR}=\dfrac{CA}{RP}$.
  \item[(c)] No. $\tr ABC\sim\tr QPR$ says that $A$ matches $Q$ and $B$ matches $P$. This is not the matching that is given.
  \item[(d)] Yes. $B$ still matches $Q$, $C$ matches $R$ and $A$ matches $P$. Only the starting vertex changed.
\end{steps}
\end{solution}
  \item In the figure, $\tr ABC$ has $AB=4$ cm, $BC=6$ cm and $CA=5$ cm. $\tr XYZ$ has $XY=12$ cm, $YZ=8$ cm and $ZX=10$ cm. Also $\angle A=\angle Z$, $\angle B=\angle Y$ and $\angle C=\angle X$.
  %FIG: two triangles side by side. Left: triangle ABC, B bottom-left, C bottom-right, A top (sides AB=4, BC=6, CA=5 drawn to scale, label the three lengths outside the sides). Right: triangle XYZ, drawn turned so that the matching is not obvious: Y at the bottom-right, Z at the bottom-left... use YZ=8, XY=12, ZX=10 to scale (twice the left triangle), X at the top. Label the three lengths outside the sides. Mark angle B and angle Y with one arc each, angle A and angle Z with two arcs each, angle C and angle X with three arcs each.
  \exfig{%
  \coordinate (A) at (1.575,2.315); \coordinate (B) at (0,0); \coordinate (C) at (4.2,0);
  \coordinate (X) at (6.575,6.945); \coordinate (Y) at (11.3,0); \coordinate (Z) at (5.7,0);
  \draw (A)--(B)--(C)--cycle;
  \draw (X)--(Y)--(Z)--cycle;
  \pic[draw,angle radius=5mm] {angle=C--B--A};
  \pic[draw,angle radius=5mm] {angle=B--A--C};
  \pic[draw,angle radius=6mm] {angle=B--A--C};
  \pic[draw,angle radius=5mm] {angle=A--C--B};
  \pic[draw,angle radius=6mm] {angle=A--C--B};
  \pic[draw,angle radius=7mm] {angle=A--C--B};
  \pic[draw,angle radius=5mm] {angle=X--Y--Z};
  \pic[draw,angle radius=5mm] {angle=Y--Z--X};
  \pic[draw,angle radius=6mm] {angle=Y--Z--X};
  \pic[draw,angle radius=5mm] {angle=Z--X--Y};
  \pic[draw,angle radius=6mm] {angle=Z--X--Y};
  \pic[draw,angle radius=7mm] {angle=Z--X--Y};
  \node[above] at (A) {$A$}; \node[below left] at (B) {$B$}; \node[below right] at (C) {$C$};
  \node[above] at (X) {$X$}; \node[below right] at (Y) {$Y$}; \node[below left] at (Z) {$Z$};
  \node at (0.35,1.45) {$4$};
  \node[below] at (2.1,-0.1) {$6$};
  \node at (3.25,1.55) {$5$};
  \node[below] at (8.5,-0.1) {$8$};
  \node at (9.45,3.85) {$12$};
  \node at (5.5,3.5) {$10$};
}
  \begin{enumerate}[label=(\alph*),itemsep=1pt]
    \item Complete the statement $\tr ABC\sim\tr\ \underline{\hspace{1.5cm}}$\,.
    \item Which side of $\tr XYZ$ corresponds to $BC$? Which side corresponds to $CA$?
    \item Find the three ratios $\dfrac{AB}{ZY}$, $\dfrac{BC}{YX}$ and $\dfrac{CA}{XZ}$. What do you notice?
  \end{enumerate}

\begin{solution}
\begin{steps}
  \item[(a)] $\tr ABC\sim\tr ZYX$ ($A$ matches $Z$, $B$ matches $Y$, $C$ matches $X$).
  \item[(b)] $BC$ corresponds to $YX$. $CA$ corresponds to $XZ$.
  \item[(c)] $\dfrac{AB}{ZY}=\dfrac{4}{8}=\dfrac12$, $\dfrac{BC}{YX}=\dfrac{6}{12}=\dfrac12$, $\dfrac{CA}{XZ}=\dfrac{5}{10}=\dfrac12$. All three ratios are equal.
\end{steps}
\end{solution}
\end{solvefirst}

\facilitatornote{Students often write the letters in alphabetical order. Ask them to start from one pair of equal angles and match the other two pairs before they write the statement.}

\rememberbox{In a similarity statement, write the vertices of the second triangle in the order that matches the first triangle. Then the sides and angles can be read off directly: $\tr ABC\sim\tr ZYX$ gives $\dfrac{AB}{ZY}=\dfrac{BC}{YX}=\dfrac{CA}{XZ}$.}

\section{Equal angles: the AAA and AA criteria}

To show that two triangles are similar, must we check all three pairs of angles \emph{and} all three ratios of sides? Solve P3 and P4 to find out.

\begin{solvefirst}
  \item On an A4 sheet, draw $BC=3$ cm. At $B$, draw a ray that makes an angle of $60^\circ$ with $BC$. At $C$, draw a ray on the same side of $BC$ that makes an angle of $40^\circ$ with $CB$. Let the two rays meet at $A$.
  Next, draw $EF=6$ cm. At $E$, draw a ray that makes an angle of $60^\circ$ with $EF$. At $F$, draw a ray on the same side that makes an angle of $40^\circ$ with $FE$. Let the two rays meet at $D$.
  \begin{enumerate}[label=(\alph*),itemsep=1pt]
    \item Find $\angle A$ and $\angle D$ without measuring. Give a reason.
    \item Measure $AB$, $CA$, $DE$ and $FD$ to the nearest millimetre.
    \item Find $\dfrac{AB}{DE}$, $\dfrac{BC}{EF}$ and $\dfrac{CA}{FD}$. What do you notice?
    \item Are the two triangles similar? Write the similarity statement.
  \end{enumerate}
\begin{solution}
\anspair{%
\begin{steps}
  \item[(a)] $\angle A=180^\circ-60^\circ-40^\circ=80^\circ$ and $\angle D=80^\circ$ (angle sum property of a triangle).
  \item[(b)] $AB\approx2.0$ cm, $CA\approx2.6$ cm, $DE\approx3.9$ cm, $FD\approx5.3$ cm (accept values within $1$ mm).
  \item[(c)] Each ratio is about $\dfrac12$, since $\dfrac{BC}{EF}=\dfrac36=\dfrac12$. The three ratios are equal.
  \item[(d)] Yes. $\tr ABC\sim\tr DEF$.
\end{steps}}{%
%FIG: small triangle ABC (BC = 3 at the bottom, angle B = 60, angle C = 40, A at the top) and next to it the larger triangle DEF (EF = 6, angle E = 60, angle F = 40, D at the top), same orientation. Mark the 60 degree angles with one arc and the 40 degree angles with two arcs. Label BC = 3 and EF = 6 below the bases.
\coordinate (A) at (0.568,0.984); \coordinate (B) at (0,0); \coordinate (C) at (1.74,0);
\coordinate (D) at (4.326,1.967); \coordinate (E) at (3.19,0); \coordinate (F) at (6.67,0);
\draw (A)--(B)--(C)--cycle;
\draw (D)--(E)--(F)--cycle;
\pic[draw,angle radius=4mm] {angle=C--B--A};
\pic[draw,angle radius=4mm] {angle=A--C--B};
\pic[draw,angle radius=4mm] {angle=F--E--D};
\pic[draw,angle radius=4mm] {angle=D--F--E};
\node[above] at (A) {$A$}; \node[left] at (B) {$B$}; \node[right] at (C) {$C$};
\node[above] at (D) {$D$}; \node[left] at (E) {$E$}; \node[right] at (F) {$F$};
\pic[draw,angle radius=5mm] {angle=A--C--B};
\pic[draw,angle radius=5mm] {angle=D--F--E};
\node[below] at ($(B)!0.5!(C)$) {$3$};
\node[below] at ($(E)!0.5!(F)$) {$6$};

}
\end{solution}
  \item A square has sides of $4$ cm. A rectangle has sides of $4$ cm and $8$ cm.
  \begin{enumerate}[label=(\alph*),itemsep=1pt]
    \item Are the angles of the square equal to the angles of the rectangle?
    \item Are the sides of the square proportional to the sides of the rectangle? Give a reason.
    \item Are the square and the rectangle similar? What does this tell you about equal angles in figures with four sides?
  \end{enumerate}

\begin{solution}
\begin{steps}
  \item[(a)] Yes. Every angle is $90^\circ$.
  \item[(b)] No. $\dfrac44=1$ but $\dfrac48=\dfrac12$.
  \item[(c)] No. In figures with four sides, equal angles do not make the sides proportional. P3 suggests that triangles behave differently.
\end{steps}
\end{solution}
\end{solvefirst}

P3 suggests that equal angles are enough for triangles. The next theorem says that this is always true.

\begin{theorem}{6.3}{AAA similarity criterion}
If in two triangles the corresponding angles are equal, then their corresponding sides are in the same ratio, and so the two triangles are similar.
\end{theorem}

%FIG: proof figure. Left: small triangle ABC (B bottom-left, C bottom-right, A top). Right: larger triangle DEF (E bottom-left, F bottom-right, D top), same shape and orientation, about 1.6 times bigger. On DE mark P with DP = AB; on DF mark Q with DQ = AC; join PQ (PQ is parallel to EF). Mark equal angles: one arc at B, E and at P (angle DPQ); two arcs at C, F and at Q (angle DQP).
\exfig{%
  \coordinate (A) at (1.8,2.7); \coordinate (B) at (0,0); \coordinate (C) at (4.5,0);
  \coordinate (D) at (8.88,4.32); \coordinate (E) at (6,0); \coordinate (F) at (13.2,0);
  \coordinate (P) at (7.08,1.62); \coordinate (Q) at (11.58,1.62);
  \draw (A)--(B)--(C)--cycle;
  \draw (D)--(E)--(F)--cycle;
  \draw (P)--(Q);
  \pic[draw,angle radius=5mm] {angle=C--B--A};
  \pic[draw,angle radius=5mm] {angle=F--E--D};
  \pic[draw,angle radius=5mm] {angle=Q--P--D};
  \pic[draw,angle radius=5mm] {angle=A--C--B};
  \pic[draw,angle radius=6mm] {angle=A--C--B};
  \pic[draw,angle radius=5mm] {angle=D--F--E};
  \pic[draw,angle radius=6mm] {angle=D--F--E};
  \pic[draw,angle radius=5mm] {angle=D--Q--P};
  \pic[draw,angle radius=6mm] {angle=D--Q--P};
  \draw[-{Stealth[length=2mm]}] ($(E)!0.45!(F)$) -- ($(E)!0.55!(F)$);
  \draw[-{Stealth[length=2mm]}] ($(P)!0.45!(Q)$) -- ($(P)!0.55!(Q)$);
  \node[above] at (A) {$A$}; \node[below left] at (B) {$B$}; \node[below right] at (C) {$C$};
  \node[above] at (D) {$D$}; \node[below left] at (E) {$E$}; \node[below right] at (F) {$F$};
  \node[left] at (6.75,1.84) {$P$};
  \node[right] at (11.86,1.90) {$Q$};
}

\textbf{Given:} $\tr ABC$ and $\tr DEF$ with $\angle A=\angle D$, $\angle B=\angle E$ and $\angle C=\angle F$.\\
\textbf{To prove:} $\dfrac{AB}{DE}=\dfrac{BC}{EF}=\dfrac{CA}{FD}$.\\
\textbf{Construction:} Mark $P$ on $DE$ such that $DP=AB$, and $Q$ on $DF$ such that $DQ=AC$. Join $PQ$.\\
\textbf{Proof:}
\begin{steps}
  \item In $\tr ABC$ and $\tr DPQ$: $AB=DP$, $\angle A=\angle D$ and $AC=DQ$. So $\tr ABC\cong\tr DPQ$ (SAS congruence).
  \item So $\angle B=\angle DPQ$ (CPCT). But $\angle B=\angle E$ (given). So $\angle DPQ=\angle E$.
  \item $\angle DPQ$ and $\angle E$ are corresponding angles made by the lines $PQ$ and $EF$ with the transversal $DE$. So $PQ\parallel EF$.
  \item In $\tr DEF$, $PQ\parallel EF$. So $\dfrac{DP}{PE}=\dfrac{DQ}{QF}$ (BPT).
  \item Turn both sides upside down and add $1$: $\dfrac{PE}{DP}+1=\dfrac{QF}{DQ}+1$, that is, $\dfrac{DE}{DP}=\dfrac{DF}{DQ}$. So $\dfrac{DP}{DE}=\dfrac{DQ}{DF}$.
  \item $DP=AB$ and $DQ=AC$. So $\dfrac{AB}{DE}=\dfrac{AC}{DF}$.
  \item In the same way, mark points on $ED$ and $EF$ with lengths $BA$ and $BC$ from $E$. The same steps give $\dfrac{AB}{DE}=\dfrac{BC}{EF}$.
\end{steps}
So $\dfrac{AB}{DE}=\dfrac{BC}{EF}=\dfrac{CA}{FD}$, and $\tr ABC\sim\tr DEF$. \hfill Hence proved.

\facilitatornote{Step 5 is the step that students find hardest. Show it with numbers first: if $\dfrac{DP}{PE}=\dfrac23$, then $\dfrac{PE}{DP}=\dfrac32$ and $\dfrac{DE}{DP}=\dfrac52$. The proof assumes $AB<DE$. If $AB=DE$, the triangles are congruent. If $AB>DE$, exchange the two triangles.}

If two angles of one triangle are equal to two angles of another triangle, the third angles are also equal, by the angle sum property. So two pairs of angles are enough.

\begin{corollary}{AA similarity criterion}
If two angles of one triangle are equal to two angles of another triangle, then the two triangles are similar.
\end{corollary}

\rememberbox{To prove two triangles similar by angles, find two pairs of equal angles. Useful reasons are: a common angle, vertically opposite angles, alternate or corresponding angles for parallel lines, and right angles.}

\begin{example}[NCERT Example 4]
In the figure, the segments $PS$ and $QR$ meet at $O$, and $PQ\parallel RS$. Prove that $\tr POQ\sim\tr SOR$.
%FIG: two segments PS and QR crossing at O, forming a "bow-tie". P top-left, Q top-right (PQ short, horizontal), R bottom-left, S bottom-right (RS longer, horizontal, parallel to PQ). O between. Arrow marks on PQ and RS to show they are parallel.
\exfig{%
  \coordinate (P) at (0,3); \coordinate (Q) at (4,3);
  \coordinate (R) at (-2,0); \coordinate (S) at (6,0);
  \coordinate (O) at (2,2);
  \draw (P)--(Q); \draw (R)--(S);
  \draw (P)--(S); \draw (Q)--(R);
  \draw[-{Stealth[length=2mm]}] ($(P)!0.45!(Q)$) -- ($(P)!0.55!(Q)$);
  \draw[-{Stealth[length=2mm]}] ($(R)!0.45!(S)$) -- ($(R)!0.55!(S)$);
  \node[above left] at (P) {$P$}; \node[above right] at (Q) {$Q$};
  \node[below left] at (R) {$R$}; \node[below right] at (S) {$S$};
  \node[left] at ($(O)+(-0.5,0)$) {$O$};
}
\textbf{Given:} $PS$ and $QR$ meet at $O$; $PQ\parallel RS$.\\
\textbf{To prove:} $\tr POQ\sim\tr SOR$.\\
\textbf{Proof:} In $\tr POQ$ and $\tr SOR$:
\begin{steps}
  \item $\angle P=\angle S$ (alternate angles, $PQ\parallel RS$, transversal $PS$).
  \item $\angle Q=\angle R$ (alternate angles, $PQ\parallel RS$, transversal $QR$).
  \item $\angle POQ=\angle SOR$ (vertically opposite angles).
\end{steps}
So $\tr POQ\sim\tr SOR$ (AAA similarity criterion). \hfill Hence proved.
\end{example}

\begin{example}
In $\tr ABC$, $\angle A=50^\circ$ and $\angle B=70^\circ$. In $\tr PQR$, $\angle P=70^\circ$ and $\angle R=60^\circ$. Also $AB=6$ cm, $BC=8$ cm and $PQ=9$ cm.
\begin{enumerate}[label=(\alph*),itemsep=1pt]
  \item Show that the two triangles are similar, and write the similarity statement.
  \item Find the length of $PR$.
\end{enumerate}
%FIG: two rough triangles side by side. Left: triangle ABC with angle A = 50, angle B = 70 marked, AB = 6 and BC = 8 marked. Right: triangle PQR with angle P = 70, angle R = 60 marked, PQ = 9 marked, PR marked "?".
\exfig{%
  \coordinate (A) at (1.436,3.947); \coordinate (B) at (0,0); \coordinate (C) at (5.6,0);
  \coordinate (P) at (6.9,0); \coordinate (Q) at (13.2,0); \coordinate (R) at (8.806,5.237);
  \draw (A)--(B)--(C)--cycle;
  \draw (P)--(Q)--(R)--cycle;
  \pic[draw,"$50^\circ$",angle radius=5mm,angle eccentricity=1.9] {angle=B--A--C};
  \pic[draw,"$70^\circ$",angle radius=5mm,angle eccentricity=1.9] {angle=C--B--A};
  \pic[draw,"$70^\circ$",angle radius=5mm,angle eccentricity=1.9] {angle=Q--P--R};
  \pic[draw,"$60^\circ$",angle radius=6mm,angle eccentricity=2] {angle=P--R--Q};
  \node[above] at (A) {$A$}; \node[below left] at (B) {$B$}; \node[below right] at (C) {$C$};
  \node[below left] at (P) {$P$}; \node[below right] at (Q) {$Q$}; \node[above] at (R) {$R$};
  \node[left] at (0.15,2.15) {$6$};
  \node[below] at (2.8,-0.1) {$8$};
  \node[below] at (10.05,-0.1) {$9$};
  \node[left] at (7.4,2.85) {$?$};
}
\textbf{Given:} $\angle A=50^\circ$, $\angle B=70^\circ$, $\angle P=70^\circ$, $\angle R=60^\circ$; $AB=6$ cm, $BC=8$ cm, $PQ=9$ cm.\\
\textbf{To find:} (a) the similarity statement; (b) $PR$.\\
\textbf{Solution:}
\begin{steps}
  \item $\angle C=180^\circ-50^\circ-70^\circ=60^\circ$ and $\angle Q=180^\circ-70^\circ-60^\circ=50^\circ$ (angle sum property).
  \item $\angle A=\angle Q=50^\circ$, $\angle B=\angle P=70^\circ$, $\angle C=\angle R=60^\circ$. So $\tr ABC\sim\tr QPR$ (AA similarity criterion).
  \item Corresponding sides are in the same ratio: $\dfrac{AB}{QP}=\dfrac{BC}{PR}$.
  \item $\dfrac{6}{9}=\dfrac{8}{PR}$, so $PR=\dfrac{8\times9}{6}=12$ cm.
\end{steps}
\textbf{Answer:} (a) $\tr ABC\sim\tr QPR$; (b) $PR=12$ cm.
\end{example}

\begin{exercise}{3A}
\item In $\tr DEF$, $\angle D=45^\circ$ and $\angle E=75^\circ$. In $\tr LMN$, $\angle M=45^\circ$ and $\angle N=60^\circ$. Are the triangles similar? If they are, write the similarity statement with the vertices in the correct order.
\begin{solution}
\textbf{Given:} $\angle D=45^\circ$, $\angle E=75^\circ$, $\angle M=45^\circ$, $\angle N=60^\circ$.\\
\textbf{To find:} whether the triangles are similar, and the similarity statement.
\begin{steps}
  \item $\angle F=180^\circ-45^\circ-75^\circ=60^\circ$ and $\angle L=180^\circ-45^\circ-60^\circ=75^\circ$ (angle sum property).
  \item $\angle D=\angle M$, $\angle E=\angle L$ and $\angle F=\angle N$.
\end{steps}
\textbf{Answer:} Yes; $\tr DEF\sim\tr MLN$ (AA similarity criterion).
\end{solution}

\item In the figure, $\tr ODC\sim\tr OBA$, $\angle BOC=125^\circ$ and $\angle CDO=70^\circ$. Find $\angle DOC$, $\angle DCO$ and $\angle OAB$.
%FIG: NCERT Fig 6.35 style. Two horizontal parallel lines. Upper line: D (left) and C (right), close together. Lower line: A (left) and B (right). Segment DB and segment CA cross at O (O nearer the top). Mark angle CDO = 70 at D (inside triangle ODC) and angle BOC = 125 at O (the angle between OB and OC, on the right side of O).
\exfig{%
  \coordinate (D) at (0,6); \coordinate (C) at (2.6,6);
  \coordinate (A) at (-1.602,0); \coordinate (B) at (2.184,0);
  \coordinate (O) at (0.889,3.557);
  \draw (D)--(C); \draw (A)--(B);
  \draw (D)--(B); \draw (C)--(A);
  \pic[draw,"$70^\circ$",angle radius=5mm,angle eccentricity=1.9] {angle=O--D--C};
  \pic[draw,"$125^\circ$",angle radius=5mm,angle eccentricity=1.9] {angle=B--O--C};
  \node[above left] at (D) {$D$}; \node[above right] at (C) {$C$};
  \node[below left] at (A) {$A$}; \node[below right] at (B) {$B$};
  \node[left] at ($(O)+(-0.5,0)$) {$O$};
}
\begin{solution}
\textbf{Given:} $\tr ODC\sim\tr OBA$, $\angle BOC=125^\circ$, $\angle CDO=70^\circ$.\\
\textbf{To find:} $\angle DOC$, $\angle DCO$, $\angle OAB$.
\begin{steps}
  \item $DOB$ is a straight line. So $\angle DOC=180^\circ-125^\circ=55^\circ$ (linear pair).
  \item In $\tr ODC$: $\angle DCO=180^\circ-70^\circ-55^\circ=55^\circ$ (angle sum property).
  \item $\tr ODC\sim\tr OBA$, so $C$ matches $A$. So $\angle OAB=\angle OCD=55^\circ$ (corresponding angles of similar triangles).
\end{steps}
\textbf{Answer:} $\angle DOC=55^\circ$, $\angle DCO=55^\circ$, $\angle OAB=55^\circ$.
\end{solution}

\item The diagonals $AC$ and $BD$ of a trapezium $ABCD$, with $AB\parallel DC$, meet at $O$. Draw a rough figure. Using a similarity criterion, show that $\dfrac{OA}{OC}=\dfrac{OB}{OD}$.
\begin{solution}
\anspair{%
\textbf{Given:} trapezium $ABCD$, $AB\parallel DC$; diagonals $AC$ and $BD$ meet at $O$.\\
\textbf{To prove:} $\dfrac{OA}{OC}=\dfrac{OB}{OD}$.\\
\textbf{Proof:} In $\tr OAB$ and $\tr OCD$:
\begin{steps}
  \item $\angle OAB=\angle OCD$ (alternate angles, $AB\parallel DC$, transversal $AC$).
  \item $\angle OBA=\angle ODC$ (alternate angles, transversal $BD$).
  \item So $\tr OAB\sim\tr OCD$ (AA similarity criterion).
  \item So $\dfrac{OA}{OC}=\dfrac{OB}{OD}$ (corresponding sides of similar triangles).
\end{steps}
Hence proved.}{%
%FIG: trapezium ABCD, AB the longer bottom side (A left, B right), DC the shorter top side (D left, C right). Diagonals AC and BD meeting at O. Arrow marks on AB and DC.
\coordinate (A) at (0,0); \coordinate (B) at (4,0);
\coordinate (D) at (0.8,2.6); \coordinate (C) at (3.2,2.6);
\coordinate (O) at (2,1.625);
\draw (A)--(B)--(C)--(D)--cycle;
\draw (A)--(C); \draw (B)--(D);
\draw[-{Stealth[length=2mm]}] ($(A)!0.45!(B)$) -- ($(A)!0.55!(B)$);
\draw[-{Stealth[length=2mm]}] ($(D)!0.45!(C)$) -- ($(D)!0.55!(C)$);
\node[below left] at (A) {$A$}; \node[below right] at (B) {$B$};
\node[above left] at (D) {$D$}; \node[above right] at (C) {$C$};
\node[right] at ($(O)+(0.15,0)$) {$O$};
}
\end{solution}

\item $S$ and $T$ are points on the sides $PR$ and $QR$ of $\tr PQR$ such that $\angle P=\angle RTS$. Draw a rough figure. Show that $\tr RPQ\sim\tr RTS$.
\begin{solution}
\anspair{%
\textbf{Given:} $S$ on $PR$, $T$ on $QR$, $\angle P=\angle RTS$.\\
\textbf{To prove:} $\tr RPQ\sim\tr RTS$.\\
\textbf{Proof:} In $\tr RPQ$ and $\tr RTS$:
\begin{steps}
  \item $\angle RPQ=\angle RTS$ (given).
  \item $\angle PRQ=\angle TRS$ (common angle).
  \item So $\tr RPQ\sim\tr RTS$ (AA similarity criterion).
\end{steps}
Hence proved.}{%
%FIG: triangle PQR with P top, Q bottom-left, R bottom-right. S on PR (nearer R), T on QR (nearer R), segment ST joined; ST is NOT parallel to PQ (it slants the other way). One arc at P and one arc at angle RTS (at T, inside triangle RTS).
\coordinate (P) at (1.5,3.6); \coordinate (Q) at (0,0); \coordinate (R) at (3.6,0);
\coordinate (S) at (2.874,1.244); \coordinate (T) at (1.933,0);
\draw (P)--(Q)--(R)--cycle;
\draw (S)--(T);
\pic[draw,angle radius=5mm] {angle=Q--P--R};
\pic[draw,angle radius=5mm] {angle=R--T--S};
\node[above] at (P) {$P$}; \node[below left] at (Q) {$Q$}; \node[below right] at (R) {$R$};
\node[right] at (S) {$S$};
\node[below] at (T) {$T$};
}
\end{solution}

\item In the figure, the altitudes $AD$ and $CE$ of $\tr ABC$ meet at $P$. Show that:
(i)~$\tr AEP\sim\tr CDP$;\quad (ii)~$\tr ABD\sim\tr CBE$;\quad (iii)~$\tr AEP\sim\tr ADB$;\quad (iv)~$\tr PDC\sim\tr BEC$.
%FIG: NCERT Fig 6.38 style. Triangle ABC with A bottom-left, B bottom-right, C top (C above a point a little right of A... C near the left). E on AB with CE perpendicular to AB (right-angle mark at E). D on BC with AD perpendicular to BC (right-angle mark at D). AD and CE meet at P.
\exfig{%
  \coordinate (A) at (0,0); \coordinate (B) at (7.2,0); \coordinate (C) at (2.2,5.4);
  \coordinate (D) at (3.876,3.590); \coordinate (E) at (2.2,0); \coordinate (P) at (2.2,2.038);
  \draw (A)--(B)--(C)--cycle;
  \draw (C)--(E); \draw (A)--(D);
  \ramark{E}{A}{C}
  \ramark{D}{C}{A}
  \node[below left] at (A) {$A$}; \node[below right] at (B) {$B$}; \node[above] at (C) {$C$};
  \node[right] at (4.23,3.91) {$D$};
  \node[below] at (E) {$E$};
  \node[below right] at ($(P)+(0.08,-0.08)$) {$P$};
}
\begin{solution}
\textbf{Given:} $AD\perp BC$, $CE\perp AB$; $AD$ and $CE$ meet at $P$.\\
\textbf{To prove:} the four similarities.
\begin{steps}
  \item[(i)] In $\tr AEP$ and $\tr CDP$: $\angle AEP=\angle CDP=90^\circ$, and $\angle APE=\angle CPD$ (vertically opposite angles). So $\tr AEP\sim\tr CDP$ (AA).
  \item[(ii)] In $\tr ABD$ and $\tr CBE$: $\angle ADB=\angle CEB=90^\circ$, and $\angle ABD=\angle CBE$ (common angle). So $\tr ABD\sim\tr CBE$ (AA).
  \item[(iii)] In $\tr AEP$ and $\tr ADB$: $\angle AEP=\angle ADB=90^\circ$, and $\angle PAE=\angle BAD$ (common angle). So $\tr AEP\sim\tr ADB$ (AA).
  \item[(iv)] In $\tr PDC$ and $\tr BEC$: $\angle PDC=\angle BEC=90^\circ$, and $\angle PCD=\angle BCE$ (common angle). So $\tr PDC\sim\tr BEC$ (AA).
\end{steps}
Hence proved.
\end{solution}

\item $E$ is a point on the side $AD$ produced of a parallelogram $ABCD$, and $BE$ meets $CD$ at $F$. Draw a rough figure. Show that $\tr ABE\sim\tr CFB$.
\begin{solution}
\anspair{%
\textbf{Given:} parallelogram $ABCD$; $E$ on $AD$ produced; $BE$ meets $CD$ at $F$.\\
\textbf{To prove:} $\tr ABE\sim\tr CFB$.\\
\textbf{Proof:} In $\tr ABE$ and $\tr CFB$:
\begin{steps}
  \item $\angle A=\angle C$ (opposite angles of a parallelogram).
  \item $\angle AEB=\angle CBF$ (alternate angles, $AE\parallel BC$, transversal $BE$).
  \item So $\tr ABE\sim\tr CFB$ (AA similarity criterion).
\end{steps}
Hence proved.}{%
%FIG: parallelogram ABCD with A bottom-left, B bottom-right, C top-right, D top-left (slanted). AD extended upward beyond D to E. Segment BE crosses DC at F.
\coordinate (A) at (0,0); \coordinate (B) at (3.4,0);
\coordinate (C) at (4.3,2); \coordinate (D) at (0.9,2);
\coordinate (E) at (1.8,4); \coordinate (F) at (2.6,2);
\draw (A)--(B)--(C)--(D)--cycle;
\draw (D)--(E);
\draw (B)--(E);
\node[below left] at (A) {$A$}; \node[below right] at (B) {$B$};
\node[above right] at (C) {$C$}; \node[left] at ($(D)+(-0.1,0)$) {$D$};
\node[above] at (E) {$E$};
\node[above right] at ($(F)+(0.05,0.05)$) {$F$};
}
\end{solution}

\item In the figure, $ABC$ and $AMP$ are two right triangles, right-angled at $B$ and $M$. Prove that:
(i)~$\tr ABC\sim\tr AMP$;\quad (ii)~$\dfrac{CA}{PA}=\dfrac{BC}{MP}$.
%FIG: NCERT Fig 6.39 style. A bottom-left, P bottom-right, B on AP between them (nearer P). C above B with CB perpendicular to AP (right-angle mark at B). Segment AC. M on AC with PM perpendicular to AC (right-angle mark at M); segment PM.
\exfig{%
  \coordinate (A) at (0,0); \coordinate (P) at (7,0); \coordinate (B) at (4.5,0);
  \coordinate (C) at (4.5,5.5); \coordinate (M) at (2.807,3.431);
  \draw (A)--(P); \draw (B)--(C);
  \draw (A)--(C);
  \draw (P)--(M);
  \ramark{B}{A}{C}
  \ramark{M}{A}{P}
  \node[below left] at (A) {$A$}; \node[below right] at (P) {$P$};
  \node[below] at (B) {$B$}; \node[above] at (C) {$C$};
  \node[above left] at (M) {$M$};
}
\begin{solution}
\textbf{Given:} $\angle ABC=90^\circ$, $\angle AMP=90^\circ$.\\
\textbf{To prove:} (i) $\tr ABC\sim\tr AMP$; (ii) $\dfrac{CA}{PA}=\dfrac{BC}{MP}$.
\begin{steps}
  \item[(i)] In $\tr ABC$ and $\tr AMP$: $\angle ABC=\angle AMP=90^\circ$, and $\angle BAC=\angle MAP$ (common angle). So $\tr ABC\sim\tr AMP$ (AA).
  \item[(ii)] Corresponding sides of similar triangles are in the same ratio: $\dfrac{CA}{PA}=\dfrac{BC}{MP}$.
\end{steps}
Hence proved.
\end{solution}

\item In the figure, $E$ is a point on the side $CB$ produced of an isosceles triangle $ABC$ with $AB=AC$. If $AD\perp BC$ and $EF\perp AC$, prove that $\tr ABD\sim\tr ECF$.
%FIG: NCERT Fig 6.40 style. Isosceles triangle ABC, A top, B bottom-left, C bottom-right. D the foot of the perpendicular from A on BC (right-angle mark). CB extended to the left beyond B to E (dashed extension). F on AC with EF perpendicular to AC (right-angle mark at F); segment EF.
\exfig{%
  \coordinate (A) at (3,5); \coordinate (B) at (0,0); \coordinate (C) at (6,0);
  \coordinate (D) at (3,0); \coordinate (E) at (-2.2,0); \coordinate (F) at (3.829,3.618);
  \draw (A)--(B)--(C)--cycle;
  \draw[dashed] (B)--(E);
  \draw (A)--(D);
  \draw (E)--(F);
  \ramark{D}{A}{C}
  \ramark{F}{E}{C}
  \node[above] at (A) {$A$}; \node[below left] at (B) {$B$}; \node[below right] at (C) {$C$};
  \node[below] at (D) {$D$}; \node[below left] at (E) {$E$};
  \node[right] at (4.13,3.80) {$F$};
}
\begin{solution}
\textbf{Given:} $AB=AC$; $E$ on $CB$ produced; $AD\perp BC$; $EF\perp AC$.\\
\textbf{To prove:} $\tr ABD\sim\tr ECF$.\\
\textbf{Proof:} In $\tr ABD$ and $\tr ECF$:
\begin{steps}
  \item $\angle ABD=\angle ACB$ (angles opposite equal sides $AC$ and $AB$), and $\angle ACB=\angle ECF$ (same angle). So $\angle ABD=\angle ECF$.
  \item $\angle ADB=\angle EFC=90^\circ$ (given).
  \item So $\tr ABD\sim\tr ECF$ (AA similarity criterion).
\end{steps}
Hence proved.
\end{solution}

\item Find the mistake. Meena is told that $\angle B=\angle Q$ and $\angle C=\angle P$ in $\tr ABC$ and $\tr PQR$. She writes, ``$\tr ABC\sim\tr PQR$, so $\dfrac{AB}{PQ}=\dfrac{BC}{QR}$.'' Explain her mistake and write the correct statement and the correct ratio.
\begin{solution}
\begin{steps}
  \item $B$ matches $Q$ and $C$ matches $P$, so $A$ matches $R$.
  \item Meena matched $A$ with $P$ and $C$ with $R$ because she wrote the letters in the same order.
\end{steps}
\textbf{Answer:} $\tr ABC\sim\tr RQP$, so $\dfrac{AB}{RQ}=\dfrac{BC}{QP}=\dfrac{CA}{PR}$.
\end{solution}
\end{exercise}
\section{Proportional sides: the SSS criterion}

Equal angles are enough to make two triangles similar. Are proportional sides also enough?

\begin{solvefirst}
  \item On an A4 sheet, construct $\tr ABC$ with $AB=3$ cm, $BC=5$ cm and $CA=6$ cm. Below it, construct $\tr DEF$ with $DE=6$ cm, $EF=10$ cm and $FD=12$ cm. Use a ruler and compass.
  \begin{enumerate}[label=(\alph*),itemsep=1pt]
    \item Find $\dfrac{AB}{DE}$, $\dfrac{BC}{EF}$ and $\dfrac{CA}{FD}$.
    \item Measure all six angles with a protractor. Which angles are equal?
    \item Are the two triangles similar? Write the similarity statement.
  \end{enumerate}
\begin{solution}
\anspair{%
\begin{steps}
  \item[(a)] $\dfrac36=\dfrac{5}{10}=\dfrac{6}{12}=\dfrac12$. The sides are proportional.
  \item[(b)] $\angle A=\angle D\approx56^\circ$, $\angle B=\angle E\approx94^\circ$, $\angle C=\angle F\approx30^\circ$ (accept $\pm1^\circ$).
  \item[(c)] Yes. $\tr ABC\sim\tr DEF$.
\end{steps}}{%
%FIG: triangle ABC with AB = 3, BC = 5, CA = 6 to scale (B bottom-left, C bottom-right, A above, angle B obtuse, about 94 degrees) and beside it triangle DEF twice as big (DE = 6, EF = 10, FD = 12), same orientation. Side lengths written outside the sides.
\coordinate (B) at (0,0); \coordinate (C) at (1.75,0); \coordinate (A) at (-0.07,1.05);
\coordinate (E) at (3.4,0); \coordinate (F) at (6.9,0); \coordinate (D) at (3.26,2.1);
\draw (A)--(B)--(C)--cycle;
\draw (D)--(E)--(F)--cycle;
\node[above] at (A) {$A$}; \node[below left] at (B) {$B$}; \node[below right] at (C) {$C$};
\node[above] at (D) {$D$}; \node[below left] at (E) {$E$}; \node[below right] at (F) {$F$};
\node[left] at ($(A)!0.5!(B)$) {$3$};
\node[below] at ($(B)!0.5!(C)$) {$5$};
\node[above right] at ($(C)!0.5!(A)$) {$6$};
\node[left] at ($(D)!0.5!(E)$) {$6$};
\node[below] at ($(E)!0.5!(F)$) {$10$};
\node[above right] at ($(F)!0.5!(D)$) {$12$};
}
\end{solution}
  \item A square has sides of $4$ cm. A rhombus has sides of $4$ cm and one angle of $60^\circ$.
  \begin{enumerate}[label=(\alph*),itemsep=1pt]
    \item Are the sides of the square proportional to the sides of the rhombus?
    \item Are the square and the rhombus similar? Give a reason.
  \end{enumerate}

\begin{solution}
\begin{steps}
  \item[(a)] Yes. Every ratio is $\dfrac44=1$.
  \item[(b)] No. The square has angles of $90^\circ$; the rhombus has angles of $60^\circ$ and $120^\circ$. In figures with four sides, proportional sides do not make the angles equal.
\end{steps}
\end{solution}
\end{solvefirst}

\begin{theorem}{6.4}{SSS similarity criterion}
If in two triangles the sides of one triangle are proportional to the sides of the other triangle, then their corresponding angles are equal, and so the two triangles are similar.
\end{theorem}

\textbf{Given:} $\tr ABC$ and $\tr DEF$ with $\dfrac{AB}{DE}=\dfrac{BC}{EF}=\dfrac{CA}{FD}$, each ratio less than $1$.\\
\textbf{To prove:} $\angle A=\angle D$, $\angle B=\angle E$ and $\angle C=\angle F$.\\
\textbf{Construction:} Mark $P$ on $DE$ such that $DP=AB$, and $Q$ on $DF$ such that $DQ=AC$. Join $PQ$ (the same figure as for Theorem 6.3).\\
\textbf{Proof:}
\begin{steps}
  \item $\dfrac{DP}{DE}=\dfrac{AB}{DE}$ and $\dfrac{DQ}{DF}=\dfrac{AC}{DF}$. These are equal (given). So $\dfrac{DP}{DE}=\dfrac{DQ}{DF}$.
  \item Turn both sides upside down and subtract $1$: $\dfrac{DE}{DP}-1=\dfrac{DF}{DQ}-1$, that is, $\dfrac{PE}{DP}=\dfrac{QF}{DQ}$. So $\dfrac{DP}{PE}=\dfrac{DQ}{QF}$.
  \item So $PQ\parallel EF$ (converse of BPT).
  \item So $\angle DPQ=\angle E$ and $\angle DQP=\angle F$ (corresponding angles). So $\tr DPQ\sim\tr DEF$ (AA similarity criterion).
  \item So $\dfrac{DP}{DE}=\dfrac{PQ}{EF}$. As $DP=AB$, this gives $\dfrac{AB}{DE}=\dfrac{PQ}{EF}$.
  \item But $\dfrac{AB}{DE}=\dfrac{BC}{EF}$ (given). So $\dfrac{PQ}{EF}=\dfrac{BC}{EF}$, and $PQ=BC$.
  \item In $\tr ABC$ and $\tr DPQ$: $AB=DP$, $AC=DQ$ and $BC=PQ$. So $\tr ABC\cong\tr DPQ$ (SSS congruence).
  \item So $\angle A=\angle D$, $\angle B=\angle DPQ=\angle E$ and $\angle C=\angle DQP=\angle F$ (CPCT and step 4).
\end{steps}
So $\tr ABC\sim\tr DEF$. \hfill Hence proved.

\rememberbox{For triangles, one condition gives the other. Equal angles make the sides proportional (Theorem 6.3), and proportional sides make the angles equal (Theorem 6.4). This is not true for figures with four or more sides (P4 and P6).}

To match the vertices when only sides are given, match the shortest side with the shortest side, the middle side with the middle side, and the longest side with the longest side. The vertex shared by two sides matches the vertex shared by the corresponding two sides.

\begin{example}[based on NCERT Example 5]
In $\tr ABC$, $AB=5$ cm, $BC=8$ cm, $CA=7$ cm, $\angle A=82^\circ$ and $\angle B=60^\circ$. In $\tr PQR$, $RQ=10$ cm, $QP=16$ cm and $PR=14$ cm. Find $\angle P$.
%FIG: two rough triangles. Left: triangle ABC with AB = 5, BC = 8, CA = 7 marked, 82 at A and 60 at B. Right: triangle PQR (oriented differently) with RQ = 10, QP = 16, PR = 14 marked, angle P marked "?".
\exfig{%
\coordinate (A) at (1.45,2.511);
\coordinate (B) at (0,0);
\coordinate (C) at (4.64,0);
\coordinate (R) at (6.6,0);
\coordinate (P) at (14.72,0);
\coordinate (Q) at (7.429,5.741);
\draw (A)--(B)--(C)--cycle;
\draw (P)--(Q)--(R)--cycle;
\pic[draw,"$82^\circ$",angle radius=5mm,angle eccentricity=1.9]{angle=B--A--C};
\pic[draw,"$60^\circ$",angle radius=5mm,angle eccentricity=1.9]{angle=C--B--A};
\pic[draw,"$?$",angle radius=5mm,angle eccentricity=1.9]{angle=Q--P--R};
\node[above=1pt] at (A) {$A$};
\node[below left=1pt] at (B) {$B$};
\node[below right=1pt] at (C) {$C$};
\node[below right=1pt] at (P) {$P$};
\node[above=1pt] at (Q) {$Q$};
\node[below left=1pt] at (R) {$R$};
\node[left=2pt] at ($(A)!0.5!(B)$) {$5$};
\node[below=2pt] at ($(B)!0.5!(C)$) {$8$};
\node[above right=2pt] at ($(C)!0.5!(A)$) {$7$};
\node[below=2pt] at ($(R)!0.5!(P)$) {$14$};
\node[left=2pt] at ($(R)!0.5!(Q)$) {$10$};
\node[above right=2pt] at ($(Q)!0.5!(P)$) {$16$};
}
\textbf{Given:} $AB=5$ cm, $BC=8$ cm, $CA=7$ cm, $\angle A=82^\circ$, $\angle B=60^\circ$; $RQ=10$ cm, $QP=16$ cm, $PR=14$ cm.\\
\textbf{To find:} $\angle P$.\\
\textbf{Solution:}
\begin{steps}
  \item Shortest with shortest, and so on: $\dfrac{AB}{RQ}=\dfrac{5}{10}=\dfrac12$, $\dfrac{BC}{QP}=\dfrac{8}{16}=\dfrac12$, $\dfrac{CA}{PR}=\dfrac{7}{14}=\dfrac12$.
  \item $AB$ matches $RQ$ and $BC$ matches $QP$, so the shared vertex $B$ matches $Q$. Then $A$ matches $R$ and $C$ matches $P$. So $\tr ABC\sim\tr RQP$ (SSS similarity criterion).
  \item So $\angle P=\angle C$ (corresponding angles of similar triangles).
  \item $\angle C=180^\circ-82^\circ-60^\circ=38^\circ$ (angle sum property).
\end{steps}
\textbf{Answer:} $\angle P=38^\circ$.
\end{example}

\begin{exercise}{3B}
\item A triangle has sides of $3$ cm, $4$ cm and $6$ cm. Which of these triangles are similar to it? Give a reason for each.\\
(a)~$6$ cm, $8$ cm, $12$ cm \quad (b)~$9$ cm, $12$ cm, $16$ cm \quad (c)~$12$ cm, $6$ cm, $8$ cm
\begin{solution}
\textbf{To find:} which triangles have sides proportional to $3$, $4$, $6$.
\begin{steps}
  \item[(a)] $\dfrac36=\dfrac48=\dfrac{6}{12}=\dfrac12$. Similar (SSS).
  \item[(b)] $\dfrac39=\dfrac{4}{12}=\dfrac13$ but $\dfrac{6}{16}=\dfrac38$. Not similar.
  \item[(c)] In order, the sides are $6$, $8$, $12$: the same as (a). Similar (SSS).
\end{steps}
\textbf{Answer:} (a) and (c).
\end{solution}

\item In $\tr ABC$, $AB=4$ cm, $BC=6$ cm and $CA=8$ cm. In $\tr DEF$, $DE=6$ cm and $FD=12$ cm. Find the length of $EF$ for which $\tr ABC\sim\tr DEF$.
\begin{solution}
\textbf{Given:} $AB=4$, $BC=6$, $CA=8$; $DE=6$, $FD=12$ (in cm).\\
\textbf{To find:} $EF$.
\begin{steps}
  \item $\dfrac{AB}{DE}=\dfrac46=\dfrac23$ and $\dfrac{CA}{FD}=\dfrac{8}{12}=\dfrac23$.
  \item For SSS we need $\dfrac{BC}{EF}=\dfrac23$, that is, $\dfrac{6}{EF}=\dfrac23$. So $EF=9$.
\end{steps}
\textbf{Answer:} $EF=9$ cm.
\end{solution}

\item Two triangles have sides of $5$ cm, $6$ cm, $7$ cm and $10$ cm, $12$ cm, $13$ cm. Ravi says, ``Two pairs of sides are in the ratio $1:2$, so the triangles are similar.'' Is Ravi correct? Explain.
\begin{solution}
\begin{steps}
  \item $\dfrac{5}{10}=\dfrac{6}{12}=\dfrac12$, but $\dfrac{7}{13}\neq\dfrac12$.
  \item SSS needs all three pairs of sides in the same ratio.
\end{steps}
\textbf{Answer:} No. The third pair is not in the ratio $1:2$, so the triangles are not similar.
\end{solution}

\item In $\tr ABC$, $AB=2$ cm, $BC=3$ cm and $CA=4$ cm. In $\tr PQR$, $PQ=6$ cm, $QR=8$ cm and $RP=4$ cm.
\begin{enumerate}[label=(\alph*),itemsep=1pt]
  \item Show that the triangles are similar and write the similarity statement.
  \item Which angle of $\tr PQR$ is equal to $\angle B$?
\end{enumerate}
\begin{solution}
\textbf{Given:} $AB=2$, $BC=3$, $CA=4$; $PQ=6$, $QR=8$, $RP=4$ (in cm).\\
\textbf{To find:} (a) the similarity statement; (b) the angle equal to $\angle B$.
\begin{steps}
  \item Shortest with shortest: $\dfrac{AB}{RP}=\dfrac24$, $\dfrac{BC}{PQ}=\dfrac36$, $\dfrac{CA}{QR}=\dfrac48$. Each is $\dfrac12$. So the triangles are similar (SSS).
  \item $AB$ matches $RP$ and $BC$ matches $PQ$, so $B$ matches $P$. Then $A$ matches $R$ and $C$ matches $Q$.
\end{steps}
\textbf{Answer:} (a) $\tr ABC\sim\tr RPQ$; (b) $\angle P$.
\end{solution}

\item A triangular park has sides of $30$ m, $40$ m and $50$ m. An architect makes a model of the park with sides of $6$ cm, $8$ cm and $10$ cm.
\begin{enumerate}[label=(\alph*),itemsep=1pt]
  \item Show that the model has the same shape as the park.
  \item In the park, the angle between the $30$ m side and the $40$ m side is $90^\circ$. Which angle of the model is $90^\circ$? Give a reason.
\end{enumerate}
\begin{solution}
\begin{steps}
  \item[(a)] $30\text{ m}=3000$ cm. $\dfrac{3000}{6}=\dfrac{4000}{8}=\dfrac{5000}{10}=500$. The sides are proportional, so the triangles are similar (SSS).
  \item[(b)] The $6$ cm and $8$ cm sides of the model match the $30$ m and $40$ m sides. Corresponding angles of similar triangles are equal.
\end{steps}
\textbf{Answer:} (a) the ratio of every pair of sides is $1:500$; (b) the angle between the $6$ cm and $8$ cm sides.
\end{solution}
\end{exercise}

\section{An angle between proportional sides: the SAS criterion}

For congruence, two sides and the angle between them are enough (SAS). Is there a similar rule for similarity?

\begin{solvefirst}
  \item On an A4 sheet, draw $\tr ABC$ with $AB=3$ cm, $\angle A=50^\circ$ and $AC=4$ cm. Draw $\tr DEF$ with $DE=6$ cm, $\angle D=50^\circ$ and $DF=8$ cm.
  \begin{enumerate}[label=(\alph*),itemsep=1pt]
    \item Find $\dfrac{AB}{DE}$ and $\dfrac{AC}{DF}$. Where is the $50^\circ$ angle in each triangle?
    \item Measure $\angle B$, $\angle C$, $\angle E$ and $\angle F$. What do you notice?
    \item Are the two triangles similar? Which criterion tells you so after (b)?
  \end{enumerate}
\begin{solution}
\begin{steps}
  \item[(a)] $\dfrac36=\dfrac48=\dfrac12$. In each triangle the $50^\circ$ angle is between the two given sides.
  \item[(b)] $\angle B=\angle E\approx82^\circ$ and $\angle C=\angle F\approx48^\circ$.
  \item[(c)] Yes, by the AA (or AAA) criterion: $\tr ABC\sim\tr DEF$.
\end{steps}
\end{solution}
  \item On an A4 sheet, draw a ray $AX$ and a ray $AY$ with $\angle XAY=30^\circ$. Mark $B$ on $AX$ with $AB=6$ cm. With centre $B$ and radius $4$ cm, draw an arc that cuts $AY$ at two points $C_1$ and $C_2$.
  \begin{enumerate}[label=(\alph*),itemsep=1pt]
    \item In $\tr ABC_1$ and $\tr ABC_2$, which two sides are equal, and which angle is equal?
    \item Is the equal angle between the two equal sides?
    \item Measure $AC_1$ and $AC_2$. Are the two triangles similar?
    \item What does this tell you about the position of the angle in a rule like SAS?
  \end{enumerate}

\begin{solution}
\anspair{%
\begin{steps}
  \item[(a)] $AB$ is common, $BC_1=BC_2=4$ cm, and $\angle A=30^\circ$ is common.
  \item[(b)] No. The angle $A$ is not between $AB$ and $BC$.
  \item[(c)] $AC_1\approx2.6$ cm and $AC_2\approx7.8$ cm. The triangles have different shapes ($\angle AC_2B$ is acute and $\angle AC_1B$ is obtuse), so they are not similar.
  \item[(d)] The equal angle must be the angle between the two pairs of proportional sides.
\end{steps}}{%
%FIG: ray AY horizontal from A to the right, ray AX at 30 degrees above it. B on AX with AB = 6. Dashed arc of the circle with centre B radius 4 cutting AY at C1 (nearer A, about 2.6 from A) and C2 (farther, about 7.8 from A). Segments BC1 and BC2. Mark 30 at A.
\coordinate (A) at (0,0);
\coordinate (B) at (3.118,1.8);
\coordinate (X) at (4.32,2.494);
\coordinate (Y) at (5.4,0);
\coordinate (C1) at (1.53,0);
\coordinate (C2) at (4.705,0);
\draw (A)--(Y);
\draw (A)--(X);
\draw (B)--(C1);
\draw (B)--(C2);
\draw[dashed] (C1) arc[start angle=-131.41,end angle=-48.59,radius=2.4];
\pic[draw,angle radius=7mm]{angle=Y--A--X};
\node[font=\tiny] at (14:0.78) {$30^\circ$};
\node[below left=1pt] at (A) {$A$};
\node[above left=1pt] at (B) {$B$};
\node[above right=1pt] at (X) {$X$};
\node[above right=1pt] at (Y) {$Y$};
\node[below left=2pt] at (C1) {$C_1$};
\node[below right=2pt] at (C2) {$C_2$};
}
\end{solution}
\end{solvefirst}

\begin{theorem}{6.5}{SAS similarity criterion}
If one angle of a triangle is equal to one angle of another triangle, and the sides including these angles are proportional, then the two triangles are similar.
\end{theorem}

\textbf{Given:} $\tr ABC$ and $\tr DEF$ with $\angle A=\angle D$ and $\dfrac{AB}{DE}=\dfrac{AC}{DF}$, each ratio less than $1$.\\
\textbf{To prove:} $\tr ABC\sim\tr DEF$.\\
\textbf{Construction:} Mark $P$ on $DE$ such that $DP=AB$, and $Q$ on $DF$ such that $DQ=AC$. Join $PQ$ (the same figure as for Theorem 6.3).\\
\textbf{Proof:}
\begin{steps}
  \item $\dfrac{DP}{DE}=\dfrac{DQ}{DF}$ (given, since $DP=AB$ and $DQ=AC$). As in the proof of Theorem 6.4, this gives $\dfrac{DP}{PE}=\dfrac{DQ}{QF}$, and so $PQ\parallel EF$ (converse of BPT).
  \item So $\angle DPQ=\angle E$ and $\angle DQP=\angle F$ (corresponding angles).
  \item In $\tr ABC$ and $\tr DPQ$: $AB=DP$, $\angle A=\angle D$ and $AC=DQ$. So $\tr ABC\cong\tr DPQ$ (SAS congruence).
  \item So $\angle B=\angle DPQ=\angle E$ and $\angle C=\angle DQP=\angle F$ (CPCT and step 2).
  \item With $\angle A=\angle D$, all three pairs of angles are equal.
\end{steps}
So $\tr ABC\sim\tr DEF$ (AAA similarity criterion). \hfill Hence proved.

\rememberbox{SAS similarity: one pair of equal angles, and the two pairs of sides that \emph{make} these angles in the same ratio. An equal angle that is not between the two sides is not enough (P8).}

\begin{example}[NCERT Example 6]
In the figure, the segments $AB$ and $CD$ meet at $O$, and $OA\cdot OB=OC\cdot OD$. Show that $\angle A=\angle C$ and $\angle B=\angle D$.
%FIG: segments AB and CD crossing at O. A top-left, B bottom-right, D bottom-left, C top-right. Join AD (left side) and CB (right side), so the triangles AOD and COB are on the left and the right of O. OA a little shorter than OC.
\exfig{%
\coordinate (O) at (0,0);
\coordinate (A) at (-3.464,2);
\coordinate (B) at (5.196,-3);
\coordinate (C) at (4.33,2.5);
\coordinate (D) at (-4.157,-2.4);
\draw (A)--(B);
\draw (C)--(D);
\draw (A)--(D);
\draw (C)--(B);
\node[above left=1pt] at (A) {$A$};
\node[below right=1pt] at (B) {$B$};
\node[above right=1pt] at (C) {$C$};
\node[below left=1pt] at (D) {$D$};
\node[right=1pt] at (0.62,0) {$O$};
}
\textbf{Given:} $AB$ and $CD$ meet at $O$; $OA\cdot OB=OC\cdot OD$.\\
\textbf{To prove:} $\angle A=\angle C$ and $\angle B=\angle D$.\\
\textbf{Proof:}
\begin{steps}
  \item $OA\cdot OB=OC\cdot OD$. Divide both sides by $OC\cdot OB$: $\dfrac{OA}{OC}=\dfrac{OD}{OB}$.
  \item $\angle AOD=\angle COB$ (vertically opposite angles).
  \item In $\tr AOD$ and $\tr COB$, the angles at $O$ are equal, and the sides including them are proportional. So $\tr AOD\sim\tr COB$ (SAS similarity criterion).
  \item So $\angle A=\angle C$ and $\angle D=\angle B$ (corresponding angles of similar triangles).
\end{steps}
Hence proved.
\end{example}

\begin{example}
In $\tr ABC$, $D$ is a point on $AB$ and $E$ is a point on $AC$, with $AD=3$ cm, $AE=4$ cm, $AB=12$ cm, $AC=9$ cm and $DE=5$ cm.
\begin{enumerate}[label=(\alph*),itemsep=1pt]
  \item Show that $\tr ADE\sim\tr ACB$.
  \item Find $BC$.
\end{enumerate}
%FIG: triangle ABC with A at the bottom-left (right angle at A, but do not mark it), AB = 12 along the bottom to the right, AC = 9 going up. D on AB with AD = 3, E on AC with AE = 4. Segment DE (it slopes the other way from BC; it is NOT parallel to BC). Mark AD = 3, AB = 12, AE = 4, AC = 9, DE = 5, BC = "?".
\exfig{%
\coordinate (A) at (0,0); \coordinate (B) at (8,0); \coordinate (C) at (0,6);
\coordinate (D) at (2,0); \coordinate (E) at (0,2.667);
\draw (A)--(B)--(C)--cycle;
\draw (D)--(E);
\node[below left] at (A) {$A$}; \node[below right] at (B) {$B$}; \node[above left] at (C) {$C$};
\node[below] at (D) {$D$}; \node[left] at (E) {$E$};
\node[below] at ($(A)!0.5!(D)$) {$3$};
\node[right] at ($(A)!0.5!(E)$) {$4$};
\node[above right] at ($(D)!0.5!(E)$) {$5$};
\node[above right] at ($(B)!0.5!(C)$) {$?$};
% whole lengths AB = 12 and AC = 9 shown with dimension lines
\draw[{Stealth[length=1.6mm]}-{Stealth[length=1.6mm]},thin] (0,-0.9) -- (8,-0.9) node[midway,fill=white,inner sep=1pt] {$12$};
\draw[{Stealth[length=1.6mm]}-{Stealth[length=1.6mm]},thin] (-1.1,0) -- (-1.1,6) node[midway,fill=white,inner sep=1pt] {$9$};
}
\textbf{Given:} $AD=3$, $AE=4$, $AB=12$, $AC=9$, $DE=5$ (in cm).\\
\textbf{To find:} (a) proof that $\tr ADE\sim\tr ACB$; (b) $BC$.\\
\textbf{Solution:}
\begin{steps}
  \item $\dfrac{AD}{AC}=\dfrac39=\dfrac13$ and $\dfrac{AE}{AB}=\dfrac{4}{12}=\dfrac13$.
  \item $\angle DAE=\angle CAB$ (common angle). This angle is between the two pairs of sides.
  \item So $\tr ADE\sim\tr ACB$ (SAS similarity criterion), with $D$ matching $C$ and $E$ matching $B$.
  \item $\dfrac{DE}{CB}=\dfrac13$, so $\dfrac{5}{BC}=\dfrac13$ and $BC=15$ cm.
\end{steps}
\textbf{Answer:} (b) $BC=15$ cm.
\end{example}

\facilitatornote{Students often write $\tr ADE\sim\tr ABC$ here. Then they use $\dfrac{AD}{AB}=\dfrac{3}{12}$ and $\dfrac{AE}{AC}=\dfrac49$, which are not equal. $DE$ is not parallel to $BC$ in this figure, so BPT does not apply. Ask: ``Which side is the shorter side next to the common angle in each triangle?''}

\begin{exercise}{3C}
\item In the figure, $\dfrac{QR}{QS}=\dfrac{QT}{PR}$ and $\angle1=\angle2$. Show that $\tr PQS\sim\tr TQR$.
%FIG: NCERT Fig 6.36 style. Q at the left on a horizontal line QR, S on QR near R, R at the right end. P above, near S (a little left of S). T above P, a little to the right (T further from Q than P, on QP produced). Segments QP and PT on one line (Q, P, T collinear), PS, PR, TR. Angle 1 at Q (angle PQR), angle 2 at R (angle PRQ), marked with labels 1 and 2.
\exfigs{.5}{%
\coordinate (Q) at (0,0);
\coordinate (R) at (10,0);
\coordinate (S) at (7,0);
\coordinate (P) at (5,7);
\coordinate (T) at (6.5,9.1);
\draw (Q)--(R);
\draw (Q)--(T);
\draw (P)--(S);
\draw (P)--(R);
\draw (T)--(R);
\node[left=2pt] at (Q) {$Q$};
\node[right=2pt] at (R) {$R$};
\node[below=2pt] at (S) {$S$};
\node[right=1pt] at (5.65,6.8) {$P$};
\node[right=2pt] at (T) {$T$};
\pic[draw,"$1$",angle radius=5mm,angle eccentricity=1.8]{angle=R--Q--P};
\pic[draw,"$2$",angle radius=5mm,angle eccentricity=1.8]{angle=P--R--Q};
}
\begin{solution}
\textbf{Given:} $\dfrac{QR}{QS}=\dfrac{QT}{PR}$; $\angle1=\angle2$.\\
\textbf{To prove:} $\tr PQS\sim\tr TQR$.\\
\textbf{Proof:}
\begin{steps}
  \item In $\tr PQR$, $\angle PQR=\angle PRQ$ ($\angle1=\angle2$). So $PR=PQ$ (sides opposite equal angles).
  \item So $\dfrac{QR}{QS}=\dfrac{QT}{PQ}$, that is, $\dfrac{QS}{QR}=\dfrac{PQ}{QT}$.
  \item In $\tr PQS$ and $\tr TQR$: $\angle PQS=\angle TQR$ (common angle $\angle1$), and $\dfrac{PQ}{TQ}=\dfrac{QS}{QR}$.
  \item So $\tr PQS\sim\tr TQR$ (SAS similarity criterion).
\end{steps}
Hence proved.
\end{solution}

\item In the figure, $D$ is on $AB$, $E$ is on $AC$, and $\tr ABE\cong\tr ACD$. Show that $\tr ADE\sim\tr ABC$.
%FIG: NCERT Fig 6.37 style. Isosceles-looking triangle ABC, A top, B bottom-left, C bottom-right. D on AB and E on AC at the same height; DE joined. BE and CD joined (crossing).
\exfigs{.5}{%
\coordinate (A) at (0,7);
\coordinate (B) at (-4,0);
\coordinate (C) at (4,0);
\coordinate (D) at (-2.4,2.8);
\coordinate (E) at (2.4,2.8);
\draw (A)--(B)--(C)--cycle;
\draw (D)--(E);
\draw (B)--(E);
\draw (C)--(D);
\node[above=2pt] at (A) {$A$};
\node[below left=1pt] at (B) {$B$};
\node[below right=1pt] at (C) {$C$};
\node[left=1pt] at (-2.4,3.15) {$D$};
\node[right=1pt] at (2.4,3.15) {$E$};
}
\begin{solution}
\textbf{Given:} $\tr ABE\cong\tr ACD$.\\
\textbf{To prove:} $\tr ADE\sim\tr ABC$.\\
\textbf{Proof:}
\begin{steps}
  \item $AB=AC$ and $AE=AD$ (CPCT).
  \item So $\dfrac{AD}{AB}=\dfrac{AE}{AC}$.
  \item $\angle DAE=\angle BAC$ (common angle).
  \item So $\tr ADE\sim\tr ABC$ (SAS similarity criterion).
\end{steps}
Hence proved.
\end{solution}

\item In $\tr ABC$ and $\tr PQR$, $\angle B=\angle Q$, $AB=6$ cm, $BC=9$ cm, $PQ=4$ cm and $QR=6$ cm. If $AC=9$ cm, find $PR$.
\begin{solution}
\anspair{%
\textbf{Given:} $\angle B=\angle Q$; $AB=6$, $BC=9$, $PQ=4$, $QR=6$, $AC=9$ (in cm).\\
\textbf{To find:} $PR$.
\begin{steps}
  \item $\dfrac{AB}{PQ}=\dfrac64=\dfrac32$ and $\dfrac{BC}{QR}=\dfrac96=\dfrac32$.
  \item The equal angles $B$ and $Q$ are between these sides. So $\tr ABC\sim\tr PQR$ (SAS).
  \item $\dfrac{AC}{PR}=\dfrac32$, so $\dfrac{9}{PR}=\dfrac32$ and $PR=6$ cm.
\end{steps}
\textbf{Answer:} $PR=6$ cm.}{%
%FIG: two rough triangles: ABC with AB = 6, BC = 9, AC = 9, angle B marked with an arc; PQR with PQ = 4, QR = 6, angle Q marked with an arc, PR = "?".
\coordinate (B) at (0,0);
\coordinate (C) at (2.97,0);
\coordinate (A) at (0.66,1.867);
\coordinate (Q) at (4.76,0);
\coordinate (R) at (6.74,0);
\coordinate (P) at (5.2,1.245);
\draw (A)--(B)--(C)--cycle;
\draw (P)--(Q)--(R)--cycle;
\pic[draw,angle radius=3mm]{angle=C--B--A};
\pic[draw,angle radius=3mm]{angle=R--Q--P};
\node[above=1pt] at (A) {$A$};
\node[below left=1pt] at (B) {$B$};
\node[below right=1pt] at (C) {$C$};
\node[above=1pt] at (P) {$P$};
\node[below left=1pt] at (Q) {$Q$};
\node[below right=1pt] at (R) {$R$};
\node[left=2pt] at ($(A)!0.5!(B)$) {$6$};
\node[below=2pt] at ($(B)!0.5!(C)$) {$9$};
\node[above right=2pt] at ($(C)!0.5!(A)$) {$9$};
\node[left=2pt] at ($(P)!0.5!(Q)$) {$4$};
\node[below=2pt] at ($(Q)!0.5!(R)$) {$6$};
\node[above right=2pt] at ($(R)!0.5!(P)$) {$?$};
}
\end{solution}

\item In $\tr ABC$ and $\tr PQR$, $\dfrac{AB}{PQ}=\dfrac{AC}{PR}$ and $\angle B=\angle Q$. Asha says that the triangles are similar by SAS. Is she correct? Explain, and refer to P8.
\begin{solution}
\begin{steps}
  \item The sides $AB$ and $AC$ meet at $A$, not at $B$. So $\angle B$ is not the angle between them.
  \item P8 shows two triangles with two equal sides and an equal angle that is not between them. They are not similar.
\end{steps}
\textbf{Answer:} No. SAS needs the equal angle to be between the proportional sides, that is, $\angle A=\angle P$.
\end{solution}

\item Two straight roads start from a junction $J$. Village $A$ is on one road, $3$ km from $J$, and village $B$ is on the other road, $5$ km from $J$. On a map, $J'A'=6$ cm, $J'B'=10$ cm, the angle between the roads is the same as on the ground, and $A'B'=8$ cm.
\begin{enumerate}[label=(\alph*),itemsep=1pt]
  \item Explain why $\tr JAB\sim\tr J'A'B'$.
  \item Find the distance $AB$ on the ground.
\end{enumerate}
\begin{solution}
\anspair{%
\textbf{Given:} $JA=3$ km, $JB=5$ km; $J'A'=6$ cm, $J'B'=10$ cm, $A'B'=8$ cm; $\angle J=\angle J'$.\\
\textbf{To find:} (a) the reason for similarity; (b) $AB$.
\begin{steps}
  \item[(a)] $\dfrac{JA}{J'A'}=\dfrac{3\text{ km}}{6\text{ cm}}$ and $\dfrac{JB}{J'B'}=\dfrac{5\text{ km}}{10\text{ cm}}$: each is $\dfrac12$ km per cm. The angles at $J$ and $J'$ are equal and lie between these sides. So the triangles are similar (SAS).
  \item[(b)] $AB=\dfrac12\times8=4$ km.
\end{steps}
\textbf{Answer:} (b) $AB=4$ km.}{%
%FIG: rough triangle JAB with J at the left, A on the upper road (JA = 3 km), B on the lower road (JB = 5 km), AB = "?" joined.
\coordinate (J) at (0,0);
\coordinate (B) at (5,0);
\coordinate (A) at (1.8,2.4);
\draw (J)--(B)--(A)--cycle;
\node[left=2pt] at (J) {$J$};
\node[below right=1pt] at (B) {$B$};
\node[above left=1pt] at (A) {$A$};
\node[left=2pt] at ($(J)!0.5!(A)$) {$3$};
\node[below=2pt] at ($(J)!0.5!(B)$) {$5$};
\node[above right=2pt] at ($(A)!0.5!(B)$) {$?$};
}
\end{solution}

\item Construct $\tr ABC$ with $AB=4$ cm, $\angle B=60^\circ$ and $BC=3$ cm. Construct $\tr PQR$ with $PQ=8$ cm, $\angle Q=60^\circ$ and $QR=6$ cm. Without measuring, explain why $\angle A=\angle P$. Then measure both angles to check.
\begin{solution}
\textbf{Steps of construction:} draw $BC=3$ cm; construct $\angle B=60^\circ$; mark $A$ on the arm with $BA=4$ cm; join $AC$. Draw $\tr PQR$ in the same way.
\begin{steps}
  \item $\dfrac{AB}{PQ}=\dfrac48=\dfrac12$ and $\dfrac{BC}{QR}=\dfrac36=\dfrac12$, and the $60^\circ$ angles lie between these sides.
  \item So $\tr ABC\sim\tr PQR$ (SAS), and $\angle A=\angle P$ (corresponding angles).
\end{steps}
\textbf{Answer:} $\angle A=\angle P\approx46^\circ$ (measured).
\end{solution}

\item State which pairs of triangles in the figure are similar. Write the similarity criterion that you used, and write each pair of similar triangles in symbolic form.
%FIG: NCERT Fig 6.34, six pairs (i)-(vi), in a 3 x 2 grid, each pair labelled (i)...(vi) below it.
% (i) triangle ABC: angle A = 60, angle B = 80, angle C = 40; triangle PQR: angle P = 60, angle Q = 80, angle R = 40.
% (ii) triangle ABC: AB = 2, BC = 2.5, CA = 3; triangle PQR: PQ = 6, QR = 4, RP = 5 (P at the top, Q bottom-left, R bottom-right).
% (iii) triangle LMP: LM = 2.7, LP = 3, MP = 2; triangle DEF: DE = 4, DF = 6, EF = 5.
% (iv) triangle MNL: MN = 2.5, ML = 5, angle M = 70 (between MN and ML); triangle PQR: PQ = 5, QR = 10, angle Q = 70 (between QP and QR).
% (v) triangle ABC: AB = 2.5, BC = 3, angle A = 80; triangle DEF: DF = 5, EF = 6, angle F = 80.
% (vi) triangle DEF: angle D = 70, angle E = 80; triangle PQR: angle Q = 80, angle R = 30.
\begin{center}
\newcommand{\ang}[2]{\pic[draw,"\scriptsize$#1$",angle radius=3mm,angle eccentricity=2.1]{angle=#2};}
\begin{tikzpicture}[scale=.47, every node/.style={font=\scriptsize}]
% (i)
\begin{scope}
\coordinate (B) at (0,0); \coordinate (C) at (4.5,0); \coordinate (A) at (0.58,3.29);
\coordinate (Q) at (6.6,0); \coordinate (R) at (11.1,0); \coordinate (P) at (7.18,3.29);
\draw (A)--(B)--(C)--cycle; \draw (P)--(Q)--(R)--cycle;
\ang{60^\circ}{B--A--C} \ang{80^\circ}{C--B--A} \ang{40^\circ}{A--C--B}
\ang{60^\circ}{Q--P--R} \ang{80^\circ}{R--Q--P} \ang{40^\circ}{P--R--Q}
\node[above] at (A) {$A$}; \node[below left] at (B) {$B$}; \node[below right] at (C) {$C$};
\node[above] at (P) {$P$}; \node[below left] at (Q) {$Q$}; \node[below right] at (R) {$R$};
\node at (5.25,-1.4) {(i)};
\end{scope}
% (ii)
\begin{scope}[xshift=14cm]
\coordinate (B) at (0,0); \coordinate (C) at (2.5,0); \coordinate (A) at (0.25,1.984);
\coordinate (Q) at (4.5,0); \coordinate (R) at (8.5,0); \coordinate (P) at (7.875,4.961);
\draw (A)--(B)--(C)--cycle; \draw (P)--(Q)--(R)--cycle;
\node[above] at (A) {$A$}; \node[below left] at (B) {$B$}; \node[below right] at (C) {$C$};
\node[above] at (P) {$P$}; \node[below left] at (Q) {$Q$}; \node[below right] at (R) {$R$};
\node[left] at ($(A)!.5!(B)$) {$2$}; \node[below] at ($(B)!.5!(C)$) {$2.5$}; \node[above right] at ($(A)!.5!(C)$) {$3$};
\node[above left] at ($(P)!.5!(Q)$) {$6$}; \node[below] at ($(Q)!.5!(R)$) {$4$}; \node[right] at ($(P)!.5!(R)$) {$5$};
\node at (4.25,-1.4) {(ii)};
\end{scope}
% (iii)
\begin{scope}[yshift=-8cm]
\coordinate (M) at (0,0); \coordinate (P) at (2,0); \coordinate (L) at (0.57,2.64);
\coordinate (E) at (4,0); \coordinate (F) at (9,0); \coordinate (D) at (4.5,3.97);
\draw (L)--(M)--(P)--cycle; \draw (D)--(E)--(F)--cycle;
\node[above] at (L) {$L$}; \node[below left] at (M) {$M$}; \node[below right] at (P) {$P$};
\node[above] at (D) {$D$}; \node[below left] at (E) {$E$}; \node[below right] at (F) {$F$};
\node[left] at ($(L)!.5!(M)$) {$2.7$}; \node[below] at ($(M)!.5!(P)$) {$2$}; \node[right] at ($(L)!.5!(P)$) {$3$};
\node[left] at ($(D)!.5!(E)$) {$4$}; \node[below] at ($(E)!.5!(F)$) {$5$}; \node[above right] at ($(D)!.5!(F)$) {$6$};
\node at (4.5,-1.4) {(iii)};
\end{scope}
% (iv)
\begin{scope}[xshift=14cm,yshift=-8cm]
\coordinate (M) at (0.855,3.214); \coordinate (N) at (0,0.865); \coordinate (L) at (4.685,0);
\coordinate (Q) at (6.5,0); \coordinate (R) at (16.5,0); \coordinate (P) at (8.21,4.7);
\draw (M)--(N)--(L)--cycle; \draw (P)--(Q)--(R)--cycle;
\ang{70^\circ}{N--M--L} \ang{70^\circ}{R--Q--P}
\node[above] at (M) {$M$}; \node[left] at (N) {$N$}; \node[below] at (L) {$L$};
\node[above] at (P) {$P$}; \node[below left] at (Q) {$Q$}; \node[below right] at (R) {$R$};
\node[left] at ($(M)!.5!(N)$) {$2.5$}; \node[above right] at ($(M)!.5!(L)$) {$5$};
\node[left] at ($(P)!.5!(Q)$) {$5$}; \node[below] at ($(Q)!.5!(R)$) {$10$};
\node at (8.25,-1.4) {(iv)};
\end{scope}
% (v)
\begin{scope}[yshift=-16.5cm]
\coordinate (B) at (0,0); \coordinate (C) at (4.8,0); \coordinate (A) at (2.84,2.82);
\coordinate (E) at (6.9,0); \coordinate (F) at (12.9,0); \coordinate (D) at (12.03,4.92);
\draw (A)--(B)--(C)--cycle; \draw (D)--(E)--(F)--cycle;
\ang{80^\circ}{B--A--C} \ang{80^\circ}{D--F--E}
\node[above] at (A) {$A$}; \node[below left] at (B) {$B$}; \node[below right] at (C) {$C$};
\node[above] at (D) {$D$}; \node[below left] at (E) {$E$}; \node[below right] at (F) {$F$};
\node[above left] at ($(A)!.5!(B)$) {$2.5$}; \node[below] at ($(B)!.5!(C)$) {$3$};
\node[right] at ($(D)!.5!(F)$) {$5$}; \node[below] at ($(E)!.5!(F)$) {$6$};
\node at (6.15,-1.4) {(v)};
\end{scope}
% (vi)
\begin{scope}[xshift=15cm,yshift=-16.5cm]
\coordinate (E) at (0,0); \coordinate (F) at (6.4,0); \coordinate (D) at (0.59,3.35);
\coordinate (Q) at (8.6,0); \coordinate (R) at (14.6,0); \coordinate (P) at (9.155,3.144);
\draw (D)--(E)--(F)--cycle; \draw (P)--(Q)--(R)--cycle;
\ang{70^\circ}{E--D--F} \ang{80^\circ}{F--E--D}
\ang{80^\circ}{R--Q--P} \ang{30^\circ}{P--R--Q}
\node[above] at (D) {$D$}; \node[below left] at (E) {$E$}; \node[below right] at (F) {$F$};
\node[above] at (P) {$P$}; \node[below left] at (Q) {$Q$}; \node[below right] at (R) {$R$};
\node at (7,-1.4) {(vi)};
\end{scope}
\end{tikzpicture}
\end{center}
\begin{solution}
\begin{steps}
  \item[(i)] All three pairs of angles are equal. $\tr ABC\sim\tr PQR$ (AAA).
  \item[(ii)] $\dfrac{AB}{QR}=\dfrac24$, $\dfrac{BC}{RP}=\dfrac{2.5}{5}$, $\dfrac{CA}{PQ}=\dfrac36$; each is $\dfrac12$. $\tr ABC\sim\tr QRP$ (SSS).
  \item[(iii)] $\dfrac{MP}{EF}=\dfrac25$ but $\dfrac{LP}{DF}=\dfrac36=\dfrac12$. Not similar.
  \item[(iv)] $\angle M=\angle Q=70^\circ$, and $\dfrac{MN}{QP}=\dfrac{2.5}{5}=\dfrac{ML}{QR}=\dfrac{5}{10}=\dfrac12$. $\tr MNL\sim\tr QPR$ (SAS).
  \item[(v)] The $80^\circ$ angle $A$ is not between $AB$ and $BC$. Not enough information; the triangles are not shown to be similar.
  \item[(vi)] $\angle F=180^\circ-70^\circ-80^\circ=30^\circ$ and $\angle P=180^\circ-80^\circ-30^\circ=70^\circ$. $\tr DEF\sim\tr PQR$ (AA).
\end{steps}
\end{solution}
\end{exercise}
\section{Heights and distances from shadows}

At the same time of day, the sun's rays fall on the ground at the same angle everywhere nearby. A vertical object and its shadow make a right angle. So an object, its shadow and the sun's ray form a right triangle, and the triangles for different objects are similar (AA). A street lamp works in the same way for a person walking past it.

\begin{example}[NCERT Example 7]
A girl of height $90$ cm walks away from the base of a lamp-post at a speed of $120$ cm per second. The lamp is $360$ cm above the ground. Find the length of her shadow after $4$ seconds.
%FIG: lamp-post AB vertical at the left (A the lamp at the top, B the base on the ground). Girl CD vertical further right (C her head, D her feet), shorter. The ground line from B through D to E, where E is the tip of the shadow. A dashed ray from A through C to E. Right-angle marks at B and D. Mark AB = 360, CD = 90, BD = 480, DE = x.
\exfig{%
\coordinate (A) at (0,6);
\coordinate (B) at (0,0);
\coordinate (C) at (8,1.5);
\coordinate (D) at (8,0);
\coordinate (E) at (10.6667,0);
\draw (A)--(B)--(E);
\draw (C)--(D);
\draw[dashed] (A)--(E);
\ramark{B}{A}{E}
\ramark{D}{C}{E}
\node[above] at (A) {$A$};
\node[below left] at (B) {$B$};
\node[above right] at (C) {$C$};
\node[below left] at (D) {$D$};
\node[below right] at (E) {$E$};
\node[left] at ($(A)!0.5!(B)$) {$360$};
\node[left] at ($(C)!0.5!(D)$) {$90$};
\node[below] at ($(B)!0.5!(D)$) {$480$};
\node[below] at ($(D)!0.5!(E)$) {$x$};
}
\textbf{Given:} lamp height $AB=360$ cm; girl's height $CD=90$ cm; speed $120$ cm/s; time $4$ s.\\
\textbf{To find:} the length of the shadow $DE$.\\
\textbf{Solution:} Let $DE=x$ cm.
\begin{steps}
  \item Distance walked: $BD=120\times4=480$ cm.
  \item In $\tr ABE$ and $\tr CDE$: $\angle B=\angle D=90^\circ$ (the lamp-post and the girl are vertical), and $\angle E$ is common. So $\tr ABE\sim\tr CDE$ (AA).
  \item $\dfrac{BE}{DE}=\dfrac{AB}{CD}$, so $\dfrac{480+x}{x}=\dfrac{360}{90}=4$.
  \item $480+x=4x$, so $3x=480$ and $x=160$.
\end{steps}
\textbf{Answer:} The shadow is $160$ cm $=1.6$ m long.
\end{example}

\section{Lengths inside similar triangles}

In similar triangles, every pair of corresponding lengths is in the same ratio as the sides. This includes medians, altitudes and angle bisectors.

\begin{example}[NCERT Example 8]
$CM$ and $RN$ are the medians of $\tr ABC$ and $\tr PQR$, and $\tr ABC\sim\tr PQR$. Prove that:
(i)~$\tr AMC\sim\tr PNR$;\quad (ii)~$\dfrac{CM}{RN}=\dfrac{AB}{PQ}$;\quad (iii)~$\tr CMB\sim\tr RNQ$.
%FIG: triangle ABC (A bottom-left, B bottom-right, C top) with M the midpoint of AB and median CM; and a larger triangle PQR (P bottom-left, Q bottom-right, R top) with N the midpoint of PQ and median RN. Tick marks: AM = MB (one tick each), PN = NQ (two ticks each).
\exfig{%
\coordinate (A) at (0,0);
\coordinate (B) at (4,0);
\coordinate (C) at (1.7,2.8);
\coordinate (M) at (2,0);
\draw (A)--(B)--(C)--cycle;
\draw (C)--(M);
\tick{A}{M}
\tick{M}{B}
\node[below left] at (A) {$A$};
\node[below right] at (B) {$B$};
\node[above] at (C) {$C$};
\node[below] at (M) {$M$};
\coordinate (P) at (5.5,0);
\coordinate (Q) at (11.5,0);
\coordinate (R) at (8.05,4.2);
\coordinate (N) at (8.5,0);
\draw (P)--(Q)--(R)--cycle;
\draw (R)--(N);
\draw[thin] ($($(P)!0.40!(N)$)!1.2mm!90:(N)$) -- ($($(P)!0.40!(N)$)!1.2mm!-90:(N)$);
\draw[thin] ($($(P)!0.50!(N)$)!1.2mm!90:(N)$) -- ($($(P)!0.50!(N)$)!1.2mm!-90:(N)$);
\draw[thin] ($($(N)!0.40!(Q)$)!1.2mm!90:(Q)$) -- ($($(N)!0.40!(Q)$)!1.2mm!-90:(Q)$);
\draw[thin] ($($(N)!0.50!(Q)$)!1.2mm!90:(Q)$) -- ($($(N)!0.50!(Q)$)!1.2mm!-90:(Q)$);
\node[below left] at (P) {$P$};
\node[below right] at (Q) {$Q$};
\node[above] at (R) {$R$};
\node[below] at (N) {$N$};
}
\textbf{Given:} $\tr ABC\sim\tr PQR$; $M$ is the midpoint of $AB$; $N$ is the midpoint of $PQ$.\\
\textbf{To prove:} (i), (ii) and (iii) above.\\
\textbf{Proof:}
\begin{steps}
  \item $\tr ABC\sim\tr PQR$, so $\dfrac{AB}{PQ}=\dfrac{BC}{QR}=\dfrac{CA}{RP}$ and $\angle A=\angle P$, $\angle B=\angle Q$.
  \item $AB=2AM$ and $PQ=2PN$ (medians). So $\dfrac{2AM}{2PN}=\dfrac{CA}{RP}$, that is, $\dfrac{AM}{PN}=\dfrac{CA}{RP}$.
  \item[(i)] In $\tr AMC$ and $\tr PNR$: $\angle A=\angle P$, and the sides including these angles are proportional (step 2). So $\tr AMC\sim\tr PNR$ (SAS).
  \item[(ii)] From (i), $\dfrac{CM}{RN}=\dfrac{CA}{RP}$. From step 1, $\dfrac{CA}{RP}=\dfrac{AB}{PQ}$. So $\dfrac{CM}{RN}=\dfrac{AB}{PQ}$.
  \item[(iii)] $\dfrac{BM}{QN}=\dfrac{AB}{PQ}=\dfrac{BC}{QR}$ (halves of $AB$ and $PQ$), and $\angle B=\angle Q$. So $\tr CMB\sim\tr RNQ$ (SAS).
\end{steps}
Hence proved.
\end{example}

\begin{exercise}{3D}
\item A vertical pole $6$ m long casts a shadow $4$ m long on the ground. At the same time, a tower casts a shadow $28$ m long. Draw a rough figure and find the height of the tower.
\begin{solution}
\anspair{%
\textbf{Given:} pole $=6$ m, its shadow $=4$ m; tower's shadow $=28$ m.\\
\textbf{To find:} the height $h$ of the tower.
\begin{steps}
  \item The pole and its shadow, and the tower and its shadow, form two right triangles. The sun's rays make the same angle with the ground. So the triangles are similar (AA).
  \item $\dfrac{h}{6}=\dfrac{28}{4}$, so $h=6\times7=42$.
\end{steps}
\textbf{Answer:} The tower is $42$ m high.}{%
%FIG: two right triangles side by side: small one with vertical side 6 (pole) and base 4 (shadow); large one with vertical side h (tower) and base 28. Dashed sun's ray as the hypotenuse of each. Right-angle marks at the bases.
\coordinate (P) at (0,0);
\coordinate (Q) at (0,0.9231);
\coordinate (S) at (0.6154,0);
\draw (P)--(Q);
\draw (P)--(S);
\draw[dashed] (Q)--(S);
\ramark{P}{Q}{S}
\node[left] at ($(P)!0.5!(Q)$) {$6$};
\node[below] at ($(P)!0.5!(S)$) {$4$};
\coordinate (F) at (1.4657,0);
\coordinate (T) at (1.4657,6.4615);
\coordinate (U) at (5.7734,0);
\draw (F)--(T);
\draw (F)--(U);
\draw[dashed] (T)--(U);
\ramark{F}{T}{U}
\node[left] at ($(F)!0.5!(T)$) {$h$};
\node[below] at ($(F)!0.5!(U)$) {$28$};
}
\end{solution}

\item A boy $150$ cm tall stands $9$ m from the base of a lamp-post. The lamp is $6$ m above the ground.
\begin{enumerate}[label=(\alph*),itemsep=1pt]
  \item Draw a rough figure and find the length of his shadow.
  \item He walks $3$ m further away from the lamp-post. Find the new length of his shadow. How much longer is it?
\end{enumerate}
\begin{solution}
\anspair{%
\textbf{Given:} lamp $=600$ cm; boy $=150$ cm; distance from the lamp-post $=900$ cm, then $1200$ cm.\\
\textbf{To find:} the shadow $x$ in each case.
\begin{steps}
  \item The two right triangles are similar (AA), as in Example 6. So $\dfrac{\text{distance}+x}{x}=\dfrac{600}{150}=4$.
  \item[(a)] $900+x=4x$, so $x=300$ cm $=3$ m.
  \item[(b)] $1200+x=4x$, so $x=400$ cm $=4$ m. It is $1$ m longer.
\end{steps}
\textbf{Answer:} (a) $3$ m; (b) $4$ m, $1$ m longer.}{%
%FIG: lamp-post 600 on the left, boy 150 at distance 900 from it, shadow x beyond the boy; dashed ray from the lamp over the boy's head to the shadow tip; right-angle marks.
\coordinate (A) at (0,2.5);
\coordinate (B) at (0,0);
\coordinate (D) at (3.75,0);
\coordinate (C) at (3.75,0.625);
\coordinate (E) at (5,0);
\draw (A)--(B);
\draw (B)--(E);
\draw (C)--(D);
\draw[dashed] (A)--(E);
\ramark{B}{A}{E}
\ramark{D}{C}{E}
\node[left] at ($(A)!0.5!(B)$) {$600$};
\node[left] at ($(C)!0.5!(D)$) {$150$};
\node[below] at ($(B)!0.5!(D)$) {$900$};
\node[below] at ($(D)!0.5!(E)$) {$x$};
}
\end{solution}

\item Priya wants to find the height of a tree. She places a small mirror flat on the ground, $8$ m from the foot of the tree. She stands $2$ m from the mirror, on the other side, and sees the top of the tree in the mirror. Her eyes are $150$ cm above the ground. At the mirror, the ray from the top of the tree and the ray to her eyes make equal angles with the ground.
\begin{enumerate}[label=(\alph*),itemsep=1pt]
  \item Draw a rough figure. Name the two similar triangles and give the reason.
  \item Find the height of the tree.
\end{enumerate}
\begin{solution}
\anspair{%
\textbf{Given:} $TF=$ tree; mirror at $M$; $FM=8$ m; $MP=2$ m; Priya's eye height $EP=150$ cm; $\angle TMF=\angle EMP$.\\
\textbf{To find:} the height $TF$.
\begin{steps}
  \item[(a)] In $\tr TFM$ and $\tr EPM$: $\angle F=\angle P=90^\circ$ and $\angle TMF=\angle EMP$ (given). So $\tr TFM\sim\tr EPM$ (AA).
  \item[(b)] $\dfrac{TF}{EP}=\dfrac{FM}{PM}$, so $\dfrac{TF}{150}=\dfrac{800}{200}=4$ and $TF=600$ cm.
\end{steps}
\textbf{Answer:} (b) The tree is $6$ m high.}{%
%FIG: ground line with F (foot of tree) at the left, M (mirror) in the middle, P (Priya's feet) at the right; FM = 8, MP = 2. Tree TF vertical at F (tall), Priya EP vertical at P (short, E her eyes). Rays T to M and M to E (a V shape at M). Right-angle marks at F and P; equal angle marks at M on both sides.
\coordinate (F) at (0,0);
\coordinate (M) at (4,0);
\coordinate (P) at (5,0);
\coordinate (T) at (0,3);
\coordinate (E) at (5,0.75);
\draw (T)--(F)--(P);
\draw (E)--(P);
\draw (T)--(M)--(E);
\ramark{F}{T}{M}
\ramark{P}{E}{M}
\pic[draw,angle radius=3mm]{angle=T--M--F};
\pic[draw,angle radius=3mm]{angle=P--M--E};
\node[below left] at (F) {$F$};
\node[below left] at (M) {$M$};
\node[below right] at (P) {$P$};
\node[above] at (T) {$T$};
\node[above] at (E) {$E$};
\node[below] at ($(F)!0.5!(M)$) {$8$};
\node[below] at ($(M)!0.5!(P)$) {$2$};
\node[right] at ($(E)!0.5!(P)$) {$150$};
}
\end{solution}

\item $AD$ and $PM$ are medians of $\tr ABC$ and $\tr PQR$, and $\tr ABC\sim\tr PQR$. Prove that $\dfrac{AB}{PQ}=\dfrac{AD}{PM}$.
\begin{solution}
\anspair{%
\textbf{Given:} $\tr ABC\sim\tr PQR$; $D$ and $M$ are the midpoints of $BC$ and $QR$.\\
\textbf{To prove:} $\dfrac{AB}{PQ}=\dfrac{AD}{PM}$.\\
\textbf{Proof:}
\begin{steps}
  \item $\dfrac{AB}{PQ}=\dfrac{BC}{QR}=\dfrac{2BD}{2QM}=\dfrac{BD}{QM}$, and $\angle B=\angle Q$ (similar triangles).
  \item So $\tr ABD\sim\tr PQM$ (SAS).
  \item So $\dfrac{AB}{PQ}=\dfrac{AD}{PM}$ (corresponding sides).
\end{steps}
Hence proved.}{%
%FIG: triangle ABC (A top, B bottom-left, C bottom-right) with median AD to the midpoint D of BC; larger triangle PQR (P top, Q bottom-left, R bottom-right) with median PM.
\coordinate (B) at (0,0);
\coordinate (C) at (1.9,0);
\coordinate (A) at (0.75,2.2);
\coordinate (D) at (0.95,0);
\draw (A)--(B)--(C)--cycle;
\draw (A)--(D);
\node[above] at (A) {$A$};
\node[below left] at (B) {$B$};
\node[below right] at (C) {$C$};
\node[below] at (D) {$D$};
\coordinate (Q) at (3.6,0);
\coordinate (R) at (6.165,0);
\coordinate (P) at (4.6125,2.97);
\coordinate (M) at (4.8825,0);
\draw (P)--(Q)--(R)--cycle;
\draw (P)--(M);
\node[above] at (P) {$P$};
\node[below left] at (Q) {$Q$};
\node[below right] at (R) {$R$};
\node[below] at (M) {$M$};
}
\end{solution}

\item $CD$ and $GH$ are the bisectors of $\angle ACB$ and $\angle EGF$, with $D$ on $AB$ and $H$ on $FE$. If $\tr ABC\sim\tr FEG$, show that:
(i)~$\dfrac{CD}{GH}=\dfrac{AC}{FG}$;\quad (ii)~$\tr DCB\sim\tr HGE$;\quad (iii)~$\tr DCA\sim\tr HGF$.
\begin{solution}
\anspair{%
\textbf{Given:} $\tr ABC\sim\tr FEG$; $CD$ bisects $\angle ACB$; $GH$ bisects $\angle EGF$.\\
\textbf{To prove:} (i), (ii) and (iii).\\
\textbf{Proof:} From $\tr ABC\sim\tr FEG$: $\angle A=\angle F$, $\angle B=\angle E$ and $\angle ACB=\angle FGE$. So $\angle ACD=\angle FGH$ and $\angle DCB=\angle HGE$ (halves of equal angles).
\begin{steps}
  \item[(iii)] In $\tr DCA$ and $\tr HGF$: $\angle A=\angle F$ and $\angle ACD=\angle FGH$. So $\tr DCA\sim\tr HGF$ (AA).
  \item[(i)] From (iii), $\dfrac{CD}{GH}=\dfrac{AC}{FG}$.
  \item[(ii)] In $\tr DCB$ and $\tr HGE$: $\angle B=\angle E$ and $\angle DCB=\angle HGE$. So $\tr DCB\sim\tr HGE$ (AA).
\end{steps}
Hence proved.}{%
%FIG: triangle ABC (A bottom-left, B bottom-right, C top) with the angle bisector CD meeting AB at D; larger triangle FEG (F bottom-left, E bottom-right, G top) with the bisector GH meeting FE at H. Equal half-angle marks at C and at G.
\coordinate (A) at (0,0);
\coordinate (B) at (2.2,0);
\coordinate (C) at (0.7,2);
\coordinate (D) at (1.0092,0);
\draw (A)--(B)--(C)--cycle;
\draw (C)--(D);
\pic[draw,angle radius=3mm]{angle=A--C--D};
\pic[draw,angle radius=3mm]{angle=D--C--B};
\node[below left] at (A) {$A$};
\node[below right] at (B) {$B$};
\node[above] at (C) {$C$};
\node[below] at (D) {$D$};
\coordinate (F) at (3.8,0);
\coordinate (E) at (6.66,0);
\coordinate (G) at (4.71,2.6);
\coordinate (H) at (5.112,0);
\draw (F)--(E)--(G)--cycle;
\draw (G)--(H);
\pic[draw,angle radius=3mm]{angle=F--G--H};
\pic[draw,angle radius=3mm]{angle=H--G--E};
\node[below left] at (F) {$F$};
\node[below right] at (E) {$E$};
\node[above] at (G) {$G$};
\node[below] at (H) {$H$};
}
\end{solution}

\item $\tr ABC\sim\tr PQR$, with $AB=6$ cm and $PQ=9$ cm. $AD$ is the altitude from $A$ to $BC$, and $PS$ is the altitude from $P$ to $QR$. $AD=4$ cm.
\begin{enumerate}[label=(\alph*),itemsep=1pt]
  \item Prove that $\tr ABD\sim\tr PQS$.
  \item Find $PS$.
\end{enumerate}
\begin{solution}
\textbf{Given:} $\tr ABC\sim\tr PQR$; $AB=6$ cm, $PQ=9$ cm; $AD\perp BC$, $PS\perp QR$; $AD=4$ cm.\\
\textbf{To find:} (a) proof; (b) $PS$.
\begin{steps}
  \item[(a)] $\angle B=\angle Q$ (similar triangles) and $\angle ADB=\angle PSQ=90^\circ$. So $\tr ABD\sim\tr PQS$ (AA).
  \item[(b)] $\dfrac{AD}{PS}=\dfrac{AB}{PQ}$, so $\dfrac{4}{PS}=\dfrac69$ and $PS=6$ cm.
\end{steps}
\textbf{Answer:} (b) $PS=6$ cm.
\end{solution}

\item $D$ is a point on the side $BC$ of $\tr ABC$ such that $\angle ADC=\angle BAC$. Draw a rough figure. Show that $CA^2=CB\cdot CD$.
\begin{solution}
\anspair{%
\textbf{Given:} $D$ on $BC$; $\angle ADC=\angle BAC$.\\
\textbf{To prove:} $CA^2=CB\cdot CD$.\\
\textbf{Proof:} In $\tr ADC$ and $\tr BAC$:
\begin{steps}
  \item $\angle ADC=\angle BAC$ (given), and $\angle ACD=\angle BCA$ (common angle).
  \item So $\tr ADC\sim\tr BAC$ (AA), with $A$ matching $B$, $D$ matching $A$ and $C$ matching $C$.
  \item So $\dfrac{CA}{CB}=\dfrac{CD}{CA}$, and $CA^2=CB\cdot CD$.
\end{steps}
Hence proved.}{%
%FIG: triangle ABC with B bottom-left, C bottom-right, A top. D on BC nearer C; segment AD. Mark angle ADC (at D, inside triangle ADC) and angle BAC (at A) with one arc each.
\coordinate (B) at (0,0);
\coordinate (C) at (5,0);
\coordinate (A) at (2.815,2.605);
\coordinate (D) at (2.688,0);
\draw (A)--(B)--(C)--cycle;
\draw (A)--(D);
\pic[draw,angle radius=4mm]{angle=C--D--A};
\pic[draw,angle radius=4mm]{angle=B--A--C};
\node[above] at (A) {$A$};
\node[below left] at (B) {$B$};
\node[below right] at (C) {$C$};
\node[below] at (D) {$D$};
}
\end{solution}

\item In the figure, the sides $AB$ and $BC$ and the median $AD$ of $\tr ABC$ are proportional to the sides $PQ$ and $QR$ and the median $PM$ of $\tr PQR$. Show that $\tr ABC\sim\tr PQR$.
%FIG: NCERT Fig 6.41 style. Triangle ABC (A top, B bottom-left, C bottom-right) with median AD (D midpoint of BC); larger triangle PQR (P top, Q bottom-left, R bottom-right) with median PM (M midpoint of QR).
\exfig{%
\coordinate (B) at (0,0);
\coordinate (C) at (4.6,0);
\coordinate (A) at (1.9,4);
\coordinate (D) at (2.3,0);
\draw (A)--(B)--(C)--cycle;
\draw (A)--(D);
\node[above] at (A) {$A$};
\node[below left] at (B) {$B$};
\node[below right] at (C) {$C$};
\node[below] at (D) {$D$};
\coordinate (Q) at (6.1,0);
\coordinate (R) at (12.54,0);
\coordinate (P) at (8.76,5.6);
\coordinate (M) at (9.32,0);
\draw (P)--(Q)--(R)--cycle;
\draw (P)--(M);
\node[above] at (P) {$P$};
\node[below left] at (Q) {$Q$};
\node[below right] at (R) {$R$};
\node[below] at (M) {$M$};
}
\begin{solution}
\textbf{Given:} $\dfrac{AB}{PQ}=\dfrac{BC}{QR}=\dfrac{AD}{PM}$; $D$ and $M$ are the midpoints of $BC$ and $QR$.\\
\textbf{To prove:} $\tr ABC\sim\tr PQR$.\\
\textbf{Proof:}
\begin{steps}
  \item $\dfrac{BC}{QR}=\dfrac{2BD}{2QM}=\dfrac{BD}{QM}$. So $\dfrac{AB}{PQ}=\dfrac{BD}{QM}=\dfrac{AD}{PM}$.
  \item So $\tr ABD\sim\tr PQM$ (SSS), and $\angle B=\angle Q$.
  \item In $\tr ABC$ and $\tr PQR$: $\angle B=\angle Q$ and $\dfrac{AB}{PQ}=\dfrac{BC}{QR}$. So $\tr ABC\sim\tr PQR$ (SAS).
\end{steps}
Hence proved.
\end{solution}

\item The sides $AB$ and $AC$ and the median $AD$ of $\tr ABC$ are proportional to the sides $PQ$ and $PR$ and the median $PM$ of $\tr PQR$. Show that $\tr ABC\sim\tr PQR$.\\
(Hint: produce $AD$ to $E$ so that $DE=AD$, and produce $PM$ to $N$ so that $MN=PM$. Join $BE$ and $QN$.)
\begin{solution}
\anspair{%
\textbf{Given:} $\dfrac{AB}{PQ}=\dfrac{AC}{PR}=\dfrac{AD}{PM}$; $AD$ and $PM$ are medians.\\
\textbf{To prove:} $\tr ABC\sim\tr PQR$.\\
\textbf{Construction:} Produce $AD$ to $E$ with $DE=AD$, and $PM$ to $N$ with $MN=PM$. Join $BE$, $CE$, $QN$, $RN$.\\
\textbf{Proof:}
\begin{steps}
  \item The diagonals of $ABEC$ bisect each other at $D$. So $ABEC$ is a parallelogram, and $BE=AC$. In the same way, $QN=PR$.
  \item $\dfrac{AB}{PQ}=\dfrac{BE}{QN}=\dfrac{AE}{PN}$, because $BE=AC$, $QN=PR$, $AE=2AD$ and $PN=2PM$.
  \item So $\tr ABE\sim\tr PQN$ (SSS), and $\angle BAE=\angle QPN$.
  \item In the same way, $\tr ACE\sim\tr PRN$ (SSS), and $\angle CAE=\angle RPN$.
  \item Adding: $\angle BAC=\angle QPR$.
  \item In $\tr ABC$ and $\tr PQR$: $\angle A=\angle P$ and $\dfrac{AB}{PQ}=\dfrac{AC}{PR}$. So $\tr ABC\sim\tr PQR$ (SAS).
\end{steps}
Hence proved.}{%
%FIG: triangle ABC (A top, B bottom-left, C bottom-right), D the midpoint of BC, AD produced below BC to E with DE = AD (dashed), BE and CE dashed (ABEC a parallelogram).
\coordinate (A) at (1,2.4);
\coordinate (B) at (0,0);
\coordinate (C) at (3.2,0);
\coordinate (D) at (1.6,0);
\coordinate (E) at (2.2,-2.4);
\draw (A)--(B)--(C)--cycle;
\draw (A)--(D);
\draw[dashed] (D)--(E);
\draw[dashed] (B)--(E)--(C);
\node[above] at (A) {$A$};
\node[below left] at (B) {$B$};
\node[below right] at (C) {$C$};
\node[below left] at (D) {$D$};
\node[below right] at (E) {$E$};
}
\end{solution}
\end{exercise}

\ifteacher
\facilitatornote{Exercise 3D, Questions 8 and 9 are the hardest questions of the class. Question 9 needs a construction. Give the hint orally to students who are stuck, and work it on the board if time is short.}
\fi

\section{Right triangles: the RHS similarity criterion}

If the hypotenuse and one side of one right triangle are proportional to the hypotenuse and one side of another right triangle, then the two triangles are similar. This is the \textbf{RHS similarity criterion}. It is useful for right triangles, but it is not one of the main theorems of this chapter.

\section{Summary}

\begin{center}
\renewcommand{\arraystretch}{1.9}
\begin{tabular}{@{}p{0.22\linewidth}p{0.72\linewidth}@{}}
\hline
\textbf{Criterion} & \textbf{What to check for $\tr ABC\sim\tr DEF$} \\
\hline
AA (AAA) & Two pairs of equal angles. \\
SSS & $\dfrac{AB}{DE}=\dfrac{BC}{EF}=\dfrac{CA}{FD}$ \\
SAS & $\angle A=\angle D$ and $\dfrac{AB}{DE}=\dfrac{AC}{DF}$ (the angle is between the sides) \\
RHS & Right angles; hypotenuses and one pair of sides proportional. \\
\hline
\end{tabular}
\end{center}

\begin{steps}
  \item Write the vertices in matching order. Match the shortest side with the shortest side, and use the shared vertex of two sides.
  \item For triangles, equal angles and proportional sides always come together. For figures with four or more sides, one does not give the other.
  \item In similar triangles, corresponding medians, altitudes and angle bisectors are in the same ratio as the sides.
  \item Heights can be found from shadows: an object and its shadow form a right triangle, and such triangles are similar at the same time of day.
\end{steps}
\end{document}
